\documentclass[12pt,a4paper]{article}
 
\usepackage[utf8]{inputenc}
\usepackage[T1]{fontenc}
\usepackage[margin=1in]{geometry}
\usepackage{amsmath,amssymb,amsfonts}
\usepackage[ruled,linesnumbered,vlined]{algorithm2e}
\SetKwFor{ForPar}{parallel for}{}{end}
\usepackage{graphicx}
\usepackage{booktabs}
\usepackage{float}
\usepackage{parskip}
\usepackage{enumitem}
\usepackage[authoryear]{natbib}
\usepackage{xcolor}
\usepackage{listings}
\usepackage{tabularx}
\newcolumntype{R}{>{\raggedleft\arraybackslash}X}
 
\lstset{
  basicstyle=\ttfamily\small,
  keywordstyle=\color{blue},
  commentstyle=\color{gray},
  stringstyle=\color{orange},
  numbers=left,
  numberstyle=\tiny,
  breaklines=true,
  frame=single,
  columns=fullflexible
}
\DeclareMathOperator*{\argmax}{max}
\DeclareMathOperator*{\argmin}{min}
 
\usepackage[colorlinks=true, linkcolor=blue, urlcolor=blue, citecolor=blue]{hyperref}
% Custom commands
\newcommand{\sbpf}{b^{\text{pf}}}
\newcommand{\MAB}{\text{MAB}}
 
% --- FRONT PAGE CONFIGURATION ---
\title{
    \textbf{Stochastic Mixed-Integer Linear Programming for Production Planning in Salmon Farming}\\
    \large ITØK391 Master Thesis \\[0.5cm]
    \textit{Department of Informatics, University of Bergen} \\[1cm]
    \includegraphics[width=0.7\textwidth]{bilder/UiB_PositivEmblem.png}
}
 
\author{
Isak Erikstad \\ 
Ulrik Haugland \\[0.5cm]
\textit{Supervisor: Dag Haugland}
}
 
\date{\today}
 
\begin{document}
 
\maketitle
\thispagestyle{empty}
\newpage
 
\newpage
\tableofcontents
\newpage
 
% --- Main Content ---
\begin{abstract}
\noindent
This thesis addresses long-term production planning in Atlantic salmon farming under biological uncertainty. Building on a deterministic MILP baseline developed for Scale Aquaculture, we formulate a four-stage SMILP over a 60-month horizon with 81 scenarios capturing temperature-driven uncertainty in growth and mortality. Since the deterministic equivalent (DE) becomes intractable at industrially relevant fleet sizes, we propose a binary Progressive Hedging matheuristic (B-PHA) that hedges only on the binary timing decisions and resolves the continuous stocking quantities through a single extensive-form stochastic LP, exploiting a structural separation between the combinatorial and recourse axes of the problem. Evaluated on instances from 4 to 70 cages on a base-model laptop, B-PHA recovers solutions within roughly five percent of the DE optimum where the comparison is possible, scales approximately linearly in fleet size, and remains tractable at scales where DE exhausts memory entirely. Under the calibrated $\pm 2$\textdegree C scenario tree, the deterministic expected-value plan is infeasible in 27 of 81 scenarios at every fleet size considered, showing that ignoring biological uncertainty is not merely suboptimal in expectation but produces plans that cannot be executed in roughly a third of futures.
\end{abstract}
 
\section{Introduction}
 
Atlantic salmon farming follows a structured biological production cycle. Juvenile salmon, termed smolt, are transferred to open sea cages at weights typically ranging from 80 to 150 grams. Over the subsequent 14 to 22 months they are grown under continuous monitoring and automated feeding regimes until reaching harvest weight. About 75\% of fish are harvested at weights between 3 and 6 kg \citep{salmonIndex}. The central planning decisions of when to stock each cage, how much smolt to deploy, and when to harvest interact in non-trivial ways. Early stocking occupies cage capacity that could otherwise hold a different cohort, while deferred harvest allows additional weight gain at the cost of tying up biomass allowances and prolonging exposure to mortality risk. For operators managing dozens of sites and hundreds of cages across a region, the combinatorial space of feasible production schedules is substantial, motivating the application of formal optimization methods to production planning.
 
Long-term planning in salmon farming is complicated by the seasonality of growth. Fish growth is temperature-dependent and varies with season and geography. Norwegian coastal waters range from around 5\textdegree C in winter to over 15\textdegree C in summer. Because producers across the country face roughly the same temperature profile, total standing biomass follows a shared seasonal curve, at its lowest around May and peaking in October \citep{mowi2025}. This synchronization creates pressure to harvest in the same autumn-winter window. Producers attempt to smooth production across the year, but biological constraints make it difficult to deviate far from the seasonal pattern.
 
Producers also face mortality risk. Mortality is subject to abrupt increases from disease or parasitic pressure. Sea lice are the dominant concern: both the infestation itself and the treatments used to control it can cause significant losses.
 
On top of the biology, the regulatory environment imposes constraints on production: each location has a maximum allowed biomass (MAB), and regions are subject to combined quotas giving a regional MAB. Short-run salmon supply is inflexible because the fish are perishable and the production cycle is long, leaving little room to adjust output in response to short-term price movements. \citet{bergfjord2009} found that salmon farmers themselves perceive future salmon prices as the single most important source of risk, though this thesis focuses on the biological uncertainties that drive growth and mortality rather than price risk.
 
In practice, production planning in the industry is often carried out using spreadsheets, rules of thumb, and the accumulated experience of site managers. Optimization-based approaches apply mathematical optimization, often formulated as mixed-integer linear programs (MILP), to find stocking and harvesting schedules that maximize net present value subject to biological and regulatory constraints. Such models can represent the combinatorial structure of the problem well. The limitation is that deterministic models assume fixed values for parameters that are uncertain, like growth rate and mortality rate. They produce a single plan optimized for a single future. When reality deviates, the plan may perform poorly, and the planner is left reacting rather than having planned for the deviation.
 
Stochastic programming offers a framework for addressing this limitation. Rather than optimizing for one assumed future, a stochastic model considers a set of possible scenarios, each representing a different realization of uncertain parameters such as growth and mortality. Decisions that must be made before the uncertainty is resolved are chosen to perform well across the full range of outcomes, while later decisions can adapt as information is revealed. When the underlying problem also involves integer decisions, such as stocking and harvesting, it leads to a stochastic mixed-integer linear programming (SMILP) problem. The resulting models are expressive enough to capture both the combinatorial planning structure and the biological uncertainty inherent in salmon production. But they are computationally demanding: the number of variables and constraints grows with the number of scenarios, and the integer requirements complicate the decomposition methods that work well for continuous stochastic problems. Efficient solution methods for large-scale SMILP problems remain an active area of research.
 
This thesis originates from a collaboration with Scale Aquaculture, who presented us with the challenge of developing a minimum 5-year production planning tool that accounts for uncertainty in biological parameters. Building on a deterministic MILP model developed in a preceding project \citep{ErikstadHaugland2025}, we formulate and implement a multistage stochastic extension that models uncertainty in fish growth and mortality across multiple scenarios. Growth uncertainty enters the model through scenario-dependent temperature profiles, and mortality through path-dependent survival rates linked to those same profiles. To handle the resulting problem size, we adapt the Progressive Hedging algorithm, applying its consensus mechanism only to the binary timing decisions and resolving the continuous stocking quantities in a subsequent extensive-form stochastic LP. The aim is to investigate what is gained by moving from deterministic to stochastic production planning, and at what computational cost.
 
\section{Salmon Farming}
\label{sec:salmonfarming}
 
Farmed salmon is the most consumed farmed marine fish globally \citep{FAO2020, mowi2025}. Norway alone accounts for approximately 50--55\% of the global production, making farmed salmon the country's second largest export category after petroleum and the dominant component of its seafood exports \citep{ssb2026}.
 
Atlantic salmon's biological life cycle shapes both its commercial value and its risk profile. The fish undergoes a three-phase life cycle: eggs hatch as fry in freshwater rivers, fry develop into smolt, and after smoltification, the physiological adaptation to saltwater, they migrate downstream to the marine environment where the majority of growth occurs. Upstream migration and sexual maturation, the third phase, is commercially undesirable as it degrades flesh quality and is suppressed in farming practice through controlled light regimes. The commercial supply chain therefore maps directly onto the first two phases: freshwater hatcheries produce smolt, which are then transferred to marine grow-out cages/units. This thesis concerns the marine phase exclusively, from sea transfer of smolt through to harvest, as it accounts for nearly all commercial value and is the stage most exposed to biological uncertainty.
 
\begin{figure}[ht]
    \centering
    \includegraphics[width=1\linewidth]{bilder/regional_map.png}
    \caption{Schematic of a representative Norwegian salmon farming region, with locations distributed along the fjord and island coast and units/cages clustered at each location.}
    \label{fig:region_map}
\end{figure}
 
Standard salmon farming consists of placing smolt in marine cages and growing them to harvest weight, typically around four to five kilograms, over a grow-out period of roughly 14 to 22 months depending on water temperature and the production strategy employed. As a cold-blooded aquatic animal, the fish does not expend energy on thermoregulation, so feed conversion is unusually favorable and the majority of caloric intake is directed into growth. The practical consequence is that water temperature is the dominant force governing how quickly fish grow: warmer water accelerates metabolic activity and feed conversion, translating directly into faster biomass accumulation. The same thermal sensitivity, however, applies to the pathogens salmon are exposed to in open-cage systems. Of particular concern is the sea louse. Sea lice are cold blooded whose generation time shortens markedly as water warms, so sustained warm periods can allow more successive generations within one season. The resulting populations spread readily across a fjord as larvae drift on currents, and they are by far the most economically damaging biological threat in open-cage salmon farming \citep{costello2006}.
 
\begin{figure}[ht]
    \centering
    \includegraphics[width=0.8\linewidth]{bilder/Ukentlige_Laksepriser_per_vekt_2013-2024.png}
    \caption{Weekly spot price per kilogram by weight class, 2013--2024. Data from the NASDAQ Salmon Index (salesprice) history~\citep{salmonIndex}.}
    \label{fig:price_timeseries}
\end{figure}
 
\begin{figure}[ht]
    \centering
    \includegraphics[width=0.8\linewidth]{bilder/Share_and_Std.png}
    \caption{Orange (left axis): traded volume share by weight class. Blue bars (right axis): standard deviation of traded prices by weight class. Data from the NASDAQ Salmon Index, 2013--2024~\citep{salmonIndex}.}
    \label{fig:price_vol_std}
\end{figure}
 
Despite its scale, salmon farming is a highly volatile business, as shown in Figure~\ref{fig:price_timeseries}. \citet{asche2019volatility} argue that the observed increase in price volatility is not primarily driven by larger biological shocks, but by a reduction in short-run supply elasticity caused by industry consolidation, capacity constraints, and more rigid harvest schedules. Because Norwegian producers face the same seasonal temperature profile, total standing biomass peaks simultaneously across the industry each autumn, concentrating harvest in the same window and depressing spot prices in the short term \citep{mowi2025}. Prices also depend on fish quality: cohorts affected by lice infestations or early sexual maturation fetch discounted prices. A second complication arises from the weight-class structure of the salmon market. As Figure~\ref{fig:price_vol_std} illustrates, extreme weight classes are traded at much lower volume and at more volatile prices. A producer that naively maximizes revenue may steer production toward weight classes that command a high \emph{average} price but are traded thinly and unpredictably, resulting in an outcome producers would consider unacceptably risky. Future salmon prices are ultimately determined by the biological outcomes and strategic harvest decisions of every producer in the market, making them volatile and difficult to model. We therefore treat price as a given parameter and focus our modelling on the biological uncertainties, which have a clearer causal structure suitable for scenario modelling.
 
Norwegian salmon production operates under a regulatory regime built around the concept of Maximum Allowable Biomass (MAB), introduced in 2005 as the central regulatory unit of aquaculture regulation \citep{hersoug2021}. A standard commercial production licence carries an MAB of 780 tonnes (945 in the northernmost counties of Troms and Finnmark), and at no point may the standing biomass associated with the licence exceed this threshold. Each marine farm site additionally carries a location-level MAB based on local biological carrying capacity, and groups of nearby sites may be subject to coordinated fallowing requirements administered through the lice-management zones (\emph{lusesoner}) \citep{mattilsynet_brakklegging2}. Violations are sanctioned through administrative fines that are decided by the specific details of the violation. \citep{akvakulturlovenSec30, prop103L2013}.
 
A second, indirect cost of non-compliance operates through the Traffic Light System introduced in 2017 \citep{NFD2017}. The Norwegian coast is divided into thirteen production zones, each assigned a green, yellow, or red colour based on estimated sea-lice-induced mortality of wild salmonids. Producers in green zones may grow their MAB allocation by 6\% in subsequent capacity rounds, those in yellow zones must hold capacity constant, and those in red zones face mandated reductions of 6\%. The right to additional MAB is valuable, with recent auction prices on the order of NOK 200{,}000 to 305{,}000 per tonne \citep{bahr2025}, so a producer's standing in these allocation rounds is a material economic asset whose value depends, among other factors, on regulatory compliance history. The opportunity cost of a violation is therefore not exhausted by the fines alone.
 
 
\section{Literature Review}
This section reviews the optimization frameworks relevant to the planning problem set out in Sections~\ref{sec:salmonfarming} and the introduction. We begin with Mixed-Integer Linear Programming (MILP), which provides the structural foundation for our model, before introducing Stochastic Programming (SP) as the framework for handling uncertainty. Their combination yields a Stochastic Mixed-Integer Linear Program (SMILP), for which solution methods are significantly less mature than for either class of problems alone. The section closes with a review of existing applications in salmon farming, tracing how the literature has progressively tackled more realistic formulations and identifying the computational bottleneck that motivates our solution approach.
 
\subsection{Mixed-Integer Linear Programming}
 
Linear programming (LP) considers optimization problems of the form
\begin{align}
\min_{x} \quad & c^\top x \\
\text{s.t.} \quad & Ax \le b, \\
& x \ge 0, \quad x \in \mathbb{R}^n,
\end{align}
where the objective function $c^\top x$ and constraints $Ax \le b$ are linear in $x$, and all decision variables are continuous and non-negative.
 
Integer Programming (IP) extends linear programming by requiring the decision variables to take integer values. A (pure) integer linear program (ILP) can be written as
\begin{align}
\min_{x} \quad & c^\top x \\
\text{s.t.} \quad & Ax \le b, \\
& x \ge 0, \quad x \in \mathbb{Z}^n,
\end{align}
with the additional restriction that decision variables are integer-valued. A mixed-integer linear program (MILP) combines the two, allowing some variables to be continuous and others integer:
\begin{align}
\min_{x,y} \quad & c^\top x + d^\top y \\
\text{s.t.} \quad & Ax + By \le b, \\
& x \ge 0, \quad y \ge 0, \\
& x \in \mathbb{R}^{n_1}, \quad y \in \mathbb{Z}^{n_2}.
\end{align}
In the salmon production planning problem considered here, the timing of stocking and harvest events is modelled with binary variables. The associated quantities, including stocking headcount, harvest weight, and biomass, are continuous. This combination yields a MILP formulation.
 
The key structural difference is therefore not in the linearity of the objective or constraints, but in the domain of the decision variables. LP problems are convex and admit polynomial-time algorithms, whereas MILP problems are generally NP-hard due to the combinatorial nature introduced by integrality constraints \citep{nemhauser1988integer}. In practice, this means that solution times can increase dramatically with problem size, and exact algorithms can exhibit exponential worst-case running time.
 
A deterministic MILP solved for a single realization of mortality and growth parameters cannot capture the value of managerial flexibility, i.e., the ability to adapt future decisions as uncertainty resolves. Incorporating biological uncertainty explicitly is therefore essential for evaluating the true value of stochastic production planning, which motivates the stochastic programming framework introduced next.
 
 
\subsection{Stochastic Programming}
\label{sec:stochastic_programming}
 
Stochastic programming (SP) extends deterministic optimization by explicitly incorporating uncertainty into the model parameters. In a two-stage stochastic linear programming framework, decisions are divided into first-stage (here-and-now) decisions made before uncertainty is revealed, and second-stage (recourse) decisions made after observing the realization of uncertain scenarios for parameters. The sequence is: the decision-maker commits to $x$ now, nature then reveals $\xi$, and the recourse decision adapts to what was revealed. A standard two-stage stochastic linear program can be written as
\begin{align}
\min_{x} \quad & c^\top x + \mathbb{E}_\xi \left[ Q(x,\xi) \right] \\
\text{s.t.} \quad & Ax \le b, \\
& x \ge 0,
\end{align}
where $\xi$ is a random vector and $Q(x,\xi)$ denotes the optimal value of the second-stage problem given first-stage decisions $x$ and realization $\xi$.
 
Stochastic programming, contrasting to a deterministic program, incorporates the value of flexibility when planning under uncertainty. This flexibility is closely related to what is referred to as real options, namely the value of being able to make improved decisions in the future once uncertainty has been partially resolved \citep{king2012modeling}. In a stochastic programming framework, this value is captured through the recourse structure, where decisions can adapt to different realized scenarios. The benefit of a stochastic solution lies in retaining the option to respond optimally to future growth and mortality outcomes.
 
A common modelling question in stochastic programming is how strictly to enforce operational constraints across scenarios. When constraints are required to hold pathwise in every scenario, the first-stage decisions must satisfy the most restrictive scenario, and the resulting plans can become effectively worst-case in character even when the objective optimises expected value. This produces feasibility regardless of which realisation occurs, but also conservatism, with the optimiser sacrificing expected performance to accommodate scenarios that may be unlikely. An alternative, widely used when the constraints model regulatory or capacity limits whose violation carries a cost but does not halt operations, is to soften selected constraints through elastic reformulation. Each violation is measured by a non-negative variable and priced in the objective at a per-unit penalty. The resulting model remains feasible in every scenario while letting the optimiser weigh constraint violations against expected profit. This is a natural fit for risk-neutral planning problems in which regulators impose fines rather than shutdowns \citep{birge2011introduction, king2012modeling}.
 
The effectiveness of stochastic programming in capturing uncertainty depends critically on how the underlying uncertainty is modelled. The computational complexity of the stochastic model increases rapidly with the number of scenarios, and a naive implementation may lead to a combinatorial explosion in the deterministic equivalent formulation \citep{birge2011introduction, shapiro2014lectures}. In multistage stochastic programming, the number of variables and constraints grows exponentially with the number of stages \citep{king2012modeling}. To address these challenges, structural restrictions and approximation methods are commonly employed to control scenario growth while preserving the key uncertainty features relevant for decision-making \citep{king2012modeling}.
 
Two important performance measures in stochastic programming are the
\emph{Expected Value of Perfect Information} (EVPI) and the \emph{Value of
the Stochastic Solution} (VSS).
 
Let $z^{SP}$ denote the optimal objective value of the stochastic program.
Let $z^{WS}$ denote the wait-and-see value, obtained by solving the
deterministic problem separately for each scenario $\xi \in \Xi$ with
perfect knowledge of $\xi$ and aggregating the scenario-optimal objectives
by their probabilities,
\begin{equation}
  z^{WS} = \sum_{\xi \in \Xi} \pi_\xi\, z^*_\xi,
\end{equation}
where $z^*_\xi$ is the optimal objective of the deterministic problem under
realisation $\xi$. Let $z^{EEV}$ denote the \emph{expected result of using
the expected-value solution}: the first-stage decisions of a deterministic
problem solved with parameters fixed at their expectations are committed,
and the resulting plan is then evaluated under the stochastic model to
obtain its expected objective.
 
Under the minimisation convention used throughout this section, the EVPI
is defined as
\begin{equation}
  \text{EVPI} = z^{SP} - z^{WS},
\end{equation}
and measures the maximum amount a decision-maker would be willing to pay
for perfect information about future uncertainty. The VSS is defined as
\begin{equation}
  \text{VSS} = z^{EEV} - z^{SP},
\end{equation}
and quantifies the benefit of solving the full stochastic model instead of
relying on a deterministic expected-value approximation. Both quantities
are non-negative by construction: the wait-and-see problem uses strictly
more information than the stochastic program, and the stochastic program
optimises over a recourse structure that the expected-value solution
ignores.
 
The salmon production model developed in this thesis is formulated as a
maximisation of expected net present value, so the signs reverse:
$\text{EVPI} = z^{WS} - z^{SP}$ and $\text{VSS} = z^{SP} - z^{EEV}$. The
maximisation form is used when reporting results.
 
Together, these measures provide a principled basis for evaluating the
trade-offs between stochastic and deterministic planning. In the context of
salmon farming, they allow us to quantify both the cost of ignoring
biological uncertainty and the potential gain from access to perfect
information, complementing the computational analysis of solution times and
scalability presented later in this thesis.
 
Writing the stochastic program out explicitly across all scenarios, in what is known as the extensive form or deterministic equivalent (DE), quickly becomes intractable in multistage settings, motivating decomposition methods \citep{birge2011introduction, shapiro2014lectures, king2012modeling}. How finely uncertainty is discretized therefore has direct consequences for computational tractability and for the reliability of measures like VSS, which depend on the scenario set used to compute them.
 
A natural way to represent discretized uncertainty is through a scenario tree. A scenario tree encodes how the world might unfold over time by discretizing the uncertain parameters into a finite set of representative realizations. It originates from a single root node representing the current state of the world, and branches at each stage as new uncertainty is revealed, with each path from root to leaf constituting one complete scenario. This structure is particularly well-suited to multistage problems, as it explicitly encodes the information available to the decision-maker at each point in time and ensures that decisions made at a given stage depend only on the history of observed realizations, a property known as non-anticipativity.
 
 
\subsubsection{Progressive Hedging}
\label{sec:ph}

The Progressive Hedging Algorithm (PHA) of \citet{rockafellar1991scenarios}
is a scenario decomposition method for stochastic programs. Rather than
solving the full deterministic equivalent, PHA decomposes the problem by
scenario and iteratively drives the individual scenario solutions toward
agreement on the non-anticipative decisions. This decomposition structure
offers a practical advantage: the full scenario tree never needs to be held
in memory simultaneously, which keeps the dimensionality much lower and
the subproblems computationally tractable even as the number of scenarios
increases.

At each iteration, every scenario subproblem is solved independently with
a modified objective that includes (i) a linear Lagrangean term penalising
deviation from the current consensus and (ii) a quadratic proximal term
that draws each scenario's solution toward the probability-weighted
average. After all subproblems are solved, the consensus is recomputed and
the Lagrangean multipliers are updated by the standard Rockafellar--Wets
rule, which adds $\rho$ times the deviation from consensus to the existing
multiplier. For convex stochastic programs, this iterative procedure is
guaranteed to converge to the optimal solution
\citep{rockafellar1991scenarios}.

\subsubsection{Progressive Hedging for Stochastic Mixed-Integer Programs}
\label{sec:ph_milp}

When integer variables are introduced, the convergence guarantee no
longer holds: the convergence proof of \citet{rockafellar1991scenarios}
relies on convexity. As a result, integer decisions can oscillate
indefinitely between iterations without settling, and finite
termination is not guaranteed. PHA applied to stochastic mixed-integer
programs (SMILPs) is therefore a heuristic, and several techniques
have been developed for making it work in practice. We describe here
the techniques most relevant to the salmon production planning
problem; the algorithmic implementation is presented in
Section~\ref{sec:algorithm}.

\paragraph{Consensus fixing.}
\citet{watson2011progressive} introduce variable fixing on consensus as
a heuristic acceleration. They observe empirically that individual
decision variables frequently converge to common values across all
scenarios in early PH iterations and rarely move thereafter. The
heuristic exploits this: once a variable takes the same integer value
across all scenarios for sufficiently many consecutive iterations, it
is fixed at that value for all subsequent iterations, removing it from
the augmented objective. The lag before fixing guards against rare
cases in which a temporarily stable value later moves under the
influence of other variables.
Variable fixing reduces solution times by shrinking the live
combinatorial space, at the cost of small reductions in solution
quality.

\paragraph{Slamming.}
Consensus fixing requires every scenario to agree on the same integer
value, and in practice a residue of binaries often refuses to settle.
\citet{watson2011progressive} introduce \emph{slamming} as a
force-fixing mechanism for this case: once the algorithm has made
sufficient progress and further natural convergence stalls, remaining
unfixed variables are committed to a chosen value and the loop is
terminated in finitely many iterations. The specific trigger and
selection criteria in their work are tuned to the cost structure of
network design; the general idea is to detect a stall and
force-fix the oscillating binaries in order to break it. This is what our implementation
takes inspiration from, with a trigger and action redesigned for our
salmon farming production planning problem as described in
Section~\ref{sec:algorithm}.


\paragraph{Two-phase decomposition.}
The idea of using integer convergence as a stopping criterion and then recovering
the continuous variables by solving the DE with the integers fixed is introduced by
\citet{lokketangen1996progressive}, who note that once the integer components of
a solution are known, the real-valued variables are obtained by solving the
deterministic equivalent with the integer values fixed, and observe that any method
for multistage stochastic LP is suitable for that solve.
\citet{crainic2009progressive} formalise this into an explicit two-phase structure.
Phase~1 runs the PH-style loop over scenario subproblems and produces a single
binary skeleton via aggregation across scenarios. Phase~2 fixes the arcs that
reached consensus in Phase~1 and re-solves the resulting restricted stochastic
network design problem with a branch-and-bound solver; arcs that did \emph{not}
reach consensus remain binary in Phase~2, so the problem is still a MIP.
Our implementation follows this two-phase pattern, but differs in one key respect:
because $q$ is the only continuous non-anticipativity variable and all binary
timing decisions are committed by Phase~1, Phase~2 reduces to a pure stochastic
LP that is solved exactly in extensive form, with non-anticipativity on $q$
enforced structurally by indexing stocking quantities by tree node rather than
by scenario.

\paragraph{Heuristic character.}
\citet{lokketangen1996progressive} were the first to treat PH as a
practical heuristic for general SMILPs, introducing the notion of 
integer convergence: terminating PH once the integer
variables reach consensus across scenarios without waiting for the
continuous variables to fully converge, and showing empirically that
the resulting solutions are of high quality despite the absence of
theoretical guarantees. Subsequent work, notably
\citet{watson2011progressive}, has confirmed that PHA with these
adaptations produces high-quality feasible solutions on large-scale
SMILPs within reasonable time limits.
This makes it a natural candidate for the salmon production planning
problem, where the deterministic equivalent becomes intractable at
scale. The B-PHA algorithm developed in this thesis builds on the
techniques described above, and Section~\ref{sec:algorithm} sets out which
are inherited directly and which are adapted for the salmon production planning problem.

\subsubsection{Applications to salmon farming}
 
\citet{langan2011} formulate the production planning problem as a two-stage stochastic linear program, keeping the model continuous by treating harvest weight as a continuous decision rather than a binary variable. Uncertainty in growth and mortality is represented through a scenario tree, and the resulting deterministic equivalent is solved directly. The key finding is that the VSS is modest given the long deterministic tail of the planning horizon and the EVPI is still positive, motivating that biological uncertainty can have real economic value in more long-term settings.
 
\citet{rynning2012} extend the framework to a three-stage formulation by introducing smolt ordering as an explicit first-stage decision prior to saltwater production. This added recourse stage substantially increases the VSS relative to the two-stage model, confirming that the value of stochastic programming grows as more decision stages are introduced. The deterministic equivalent remains tractable at the relatively small scales considered, and a key behavioral insight is that the stochastic model distributes smolt deliveries more conservatively than the deterministic benchmark to remain feasible across scenarios.
 
\citet{haereid2012} present a multistage stochastic programming formulation for the long-term salmon production planning problem, modeled as a deterministic equivalent with non-anticipativity constraints. The paper focuses on model description rather than computational results, with the authors explicitly noting that analysis on real data is left for future work. Nonetheless, it provides a detailed blueprint for how to model the key structural features of the problem, including biomass dynamics, MAB constraints, smolt ordering, and uncertainty in growth, survival, prices, and demand, all within a multistage recourse framework spanning multiple production cycles.
 
\citet{naess2019} propose a two-stage stochastic MILP for salmon production planning over a five-year horizon. In the first stage, the total volume of smolts to be stocked is committed before uncertainty is resolved. In the second stage, after uncertainty is revealed, the model optimizes stocking distribution and harvest decisions across sites. However, the combination of integrality constraints and a growing scenario set renders the deterministic equivalent computationally intractable at full scale. The authors' 16-site, 20-scenario instance exceeds available computational resources, forcing a reduction to three representative sites and the application of Sample Average Approximation. They explicitly identify scenario decomposition as a natural direction for future work.
 
\citet{lien2022} further extend the problem by introducing early sexual maturation as an additional source of biological uncertainty in a two-stage stochastic MILP. Even with a branch-and-price approach based on Dantzig-Wolfe decomposition, the solver leaves a 7.04\% optimality gap within a two-day time limit, exploring only three branch-and-bound nodes in that period. The bottleneck appears to be the per-node cost of solving all scenario subproblems, which grows prohibitively as the number of sites and scenarios increases. This directly motivates the need for a decomposition strategy that does not rely on a global branch-and-bound tree.
 
A common thread across the larger problem formulations is that solving the stochastic problem via the deterministic equivalent becomes increasingly impractical as the number of sites and scenarios grows. \citet{naess2019} explicitly identifies scenario decomposition as a natural direction for future work, and the computational difficulties reported by \citet{lien2022} point in the same direction. The Progressive Hedging Algorithm, adopted in the present thesis, is a scenario decomposition method well-suited to this setting.
 
 
%============================================================
% CHAPTER 4: PROBLEM FORMULATION AND SOLUTION METHOD
%============================================================
\section{Problem Formulation and Solution Method}
 
The problem considered in this thesis is to determine when and how much smolt should be stocked in each cage, and correspondingly when these cohorts should be harvested, over a five-year planning horizon spanning approximately three biological production cycles. A central objective is to incorporate biological uncertainty in growth and mortality rates, and to evaluate the operational and economic implications of doing so relative to a deterministic formulation. The model extends the deterministic MILP baseline developed in the preceding project \citep{ErikstadHaugland2025}.
 
%------------------------------------------------------------
\subsection{Problem Statement}
%------------------------------------------------------------
 
Let $t = 0, \ldots, T - 1$ denote monthly decision periods over a five-year ($T=60$) planning horizon, and let $\xi$ represent stochastic biological parameters governing growth and mortality. For each period, the decision-maker must determine stocking and harvesting decisions, represented by binary variables, together with continuous headcount.
 
The objective is to maximize expected net present value over the planning horizon:
\begin{equation}
  \max_{x \in X}\; \mathbb{E}_{\xi}\bigl[f(x, \xi)\bigr]
\end{equation}
where $x$ collects all decision variables, and $X$ is the set of all $x$ satisfying biological growth dynamics, capacity constraints, regulatory biomass limits, and non-anticipativity constraints. The remainder of this chapter develops each component of this formulation in turn, moving from the objective through the biological, operational, and regulatory structure of the model, before addressing the representation of uncertainty and the choice of solution algorithm.
 
 
%------------------------------------------------------------
\subsection{Biological Modelling}
%------------------------------------------------------------
\label{sec:biology}
 
Temperature is the sole stochastic parameter in the model. Once a temperature scenario $\xi$ is drawn, both the growth trajectory and the survival rates of each cohort are fully determined. This section describes how each is modelled.
 
\noindent \textbf{Temperature} \\
Sea temperature is the central environmental driver in the model. It affects the growth rates through the effective temperature. To represent this variability, the planning horizon is divided into four consecutive stages of 15 months each. In each stage the temperature realisation takes one of three labels, \emph{cold}, \emph{normal}, or \emph{warm}, drawn from a discrete distribution, yielding $3^4 = 81$ equiprobable leaf scenarios in the scenario tree. The temperature label determines a monthly sea temperature profile in degrees Celsius, which in turn drives both the growth function and the survival regime for all cohorts during that stage. The three profiles, expressed as monthly temperatures (\textdegree C) for a representative year, are given in Table~\ref{tab:temp_profiles}.
\begin{table}[ht]
\centering
\small
\begin{tabular}{l*{12}{c}}
\toprule
 & Jan & Feb & Mar & Apr & May & Jun & Jul & Aug & Sep & Oct & Nov & Dec \\
\midrule
Cold   & 3.0 & 3.0 & 3.0 & 4.0 & 7.0  & 10.0 & 12.0 & 14.5 & 13.5 & 11.0 & 8.0  & 5.5 \\
Normal & 5.0 & 5.0 & 5.0 & 6.0 & 9.0  & 12.0 & 14.0 & 16.5 & 15.5 & 13.0 & 10.0 & 7.5 \\
Warm   & 7.0 & 7.0 & 7.0 & 8.0 & 11.0 & 14.0 & 16.0 & 18.5 & 17.5 & 15.0 & 12.0 & 9.5 \\
\bottomrule
\end{tabular}
\caption{Monthly sea temperature profiles (\textdegree C) for each temperature label. The normal profile corresponds to the observed average seasonal cycle from Norwegian coastal data provided by Scale Aquaculture. Cold and warm profiles are shifted uniformly by $\pm 2$\textdegree C.}
\label{tab:temp_profiles}
\end{table}
 
\newpage
\noindent \textbf{Growth}\\
We precompute the weight trajectory for each cohort configuration. Let $U$ be the set of all units (also denoted as \textbf{cages}). A cohort $(u,s)$ is a batch of fish in unit $u \in U$ stocked in month $s \in \mathcal{T}$. Throughout, $s \in \mathcal{T}$ denotes the stocking month of a cohort and $t \in \mathcal{T}$ denotes the current month. For each cohort and scenario, the precomputation yields a deterministic monthly weight trajectory. Weight per fish at month $t+1$, denoted $\mathrm{wpf}^{\xi}_{s,t+1}$, evolves according to the growth function
\begin{equation}
\label{eq:growth}
  \mathrm{wpf}^{\xi}_{s,t+1} = \mathrm{wpf}^{\xi}_{s,t}
    + \Delta t \cdot \alpha \cdot T_{\mathrm{eff}}(\mathrm{Temp}^{\xi}_t)
      \cdot (\mathrm{wpf}^{\xi}_{s,t})^{\beta},
\end{equation}
where $\Delta t = 30$ days per month, $\alpha$ is a growth coefficient, and $\beta$ is the physiological exponent. If raw temperatures were used directly in Equation~\eqref{eq:growth}, the model would predict monotonically increasing growth with temperature, which is biologically unrealistic. We therefore introduce the function $T_{\mathrm{eff}}$, which maps real sea temperature to an \textit{effective} growth temperature, accounting for the biological reality that Atlantic salmon growth follows a bell-shaped thermal response: it peaks around 13--14\textdegree C and declines at higher temperatures due to metabolic stress \citep{Handeland2008}.
 
\noindent \textbf{Effective growth temperature} \\
We introduce a piecewise linear effective temperature function:
\begin{equation}
\label{eq:teff}
  T_{\mathrm{eff}}(T) =
  \begin{cases}
    T & \text{if } T \leq T_{\mathrm{opt}}, \\[4pt]
    T_{\mathrm{opt}} \cdot \dfrac{T_{\max} - T}{T_{\max} - T_{\mathrm{opt}}}
      & \text{if } T_{\mathrm{opt}} < T < T_{\max}, \\[4pt]
    0 & \text{if } T \geq T_{\max},
  \end{cases}
\end{equation}
where $T_{\mathrm{opt}}$ is the temperature at which growth peaks and $T_{\max}$ is the temperature above which the model assumes growth ceases. The piecewise linear form is chosen because both parameters have direct biological interpretations and are straightforward to vary across scenarios.
 
The parameters $\alpha$, $\beta$, $T_{\mathrm{opt}}$, and $T_{\max}$ were calibrated against two sources. The first is the Thermal Growth Coefficient (TGC) model of \citet{Iwama1981}, the standard industry growth model, which updates weight daily as $w_{d+1}^{1/3} = w_d^{1/3} + (\mathrm{TGC}/1000) \cdot T_d$, with TGC values for Norwegian farmed Atlantic salmon ranging from 2.4 to 3.0 \citep{Thorarensen2011}. The second is the experimental data of \citet{Handeland2008}, who measured specific growth rates for Atlantic salmon post-smolts at four constant temperatures: 0.78\%/day at 6\textdegree C, 1.35\%/day at 10\textdegree C, 1.53\%/day at 14\textdegree C, and 1.29\%/day at 18\textdegree C, directly demonstrating the growth decline above 14\textdegree C. The resulting parameter values are
\[
  \alpha = 0.011,\quad \beta = 0.65,\quad
  T_{\mathrm{opt}} = 13.6\text{\textdegree C},\quad
  T_{\max} = 23.9\text{\textdegree C}.
\]
Figure~\ref{fig:teff} shows the fitted $T_{\mathrm{eff}}$ curve against the literature data, along with the resulting monthly effective temperature under the Norwegian profile.
 
\begin{figure}[ht]
    \centering
    \includegraphics[width=1\linewidth]{bilder/temperature_cal.png}
  \caption{Effective temperature function. (a) The piecewise $T_{\mathrm{eff}}$ curve with $T_{\mathrm{opt}} = 13.6$\textdegree C and $T_{\max} = 23.9$\textdegree C, plotted against the relative growth rates measured by \citet{Handeland2008} at four constant temperatures. The dashed line shows unchanged temperature, which assumes growth increases monotonically. The shaded region indicates the Norwegian coastal temperature range. (b) Monthly effective temperature under the normal Norwegian profile. The correction is negligible below $T_{\mathrm{opt}}$ and most pronounced in August, where 16.5\textdegree C is dampened to approximately 10.5\textdegree C.}
  \label{fig:teff}
\end{figure}


\noindent \textbf{Mortality} \\
Survival follows two monthly regimes: a baseline factor $\text{Survival rate}_{\text{normal}} = (1-0.0002)^{30} \approx 0.994$, corresponding to the industry-standard daily rate of 0.2\textperthousand, and an elevated factor $\text{Survival rate}_{\text{high}} = (1-0.001)^{30} \approx 0.970$, corresponding to 1\textperthousand\ per day. Elevated mortality is triggered only when two consecutive stages both realise warm temperatures. This represents the conditions under which sea lice can proliferate, which is the modelled cause of elevated mortality.
 
The cumulative survival from stocking period $s$ to period $t$ under scenario $\xi$ is then:
\begin{equation}
    \mathrm{s}^{\xi}_{s,t} = \prod_{\tau = s}^{t} S^{\xi}_{\tau},
\end{equation}
where $S^{\xi}_{\tau}$ is the monthly survival factor in period $\tau$ under scenario $\xi$, determined by the mortality regime described above.
 
\noindent \textbf{Survival-Adjusted Biomass trajectory}
\label{par:biomasstrajectory}\\
Given a scenario $\xi$ that fixes the temperature label for each stage, both the monthly survival factor $S^\xi_t$ and the weight per fish $\mathrm{wpf}^\xi_{s,t}$ are determined for every cohort stocked at $s$ and every period $t$. The survival-adjusted biomass per fish trajectory combines the growth function~\eqref{eq:growth} with cumulative survival:
\begin{equation}
  \label{eq:sbpf}
  \mathrm{sbpf}^{\xi}_{s,t}
    = \frac{\mathrm{wpf}^{\xi}_{s,t}}{1000} \cdot
      \mathrm{s}^{\xi}_{s,t},
\end{equation}
where the division by 1000 converts grams to kilograms. Since both the growth function and the survival factors are fully determined once the temperature path is drawn, $\mathrm{sbpf}^{\xi}_{s,t}$ is precomputed for every $(s, t, \xi)$ before the optimisation is solved. The joint effect of temperature on growth and mortality thus enters the model as a single scenario-indexed parameter rather than as additional constraints or decision variables. Figure~\ref{fig:biomass_trajectory} shows an example of how the biomass of a cohort progresses under our biological modelling.
 
Note that $\mathrm{sbpf}$ is expressed per fish of the \emph{initial} cohort size at stocking, not per fish currently alive at each point in time. This convention simplifies the formulation: the total standing biomass becomes a function of the stocking quantity and the alive indicator alone, without requiring a separate fish-count state variable tracked through time.
 
\begin{figure}[ht]
    \centering
    \includegraphics[width=1\linewidth]{bilder/bio_traj.png}
    \caption{Standing biomass over 12 months for a cohort of 50{,}000 smolt at 115\,g released in January, under daily mortality of 0.02\%. Solid lines show our growth model with the piecewise $T_{\mathrm{eff}}$ correction under the three temperature scenarios in Table~\ref{tab:temp_profiles}; dashed lines show the corresponding TGC~2.7 reference trajectories. The dotted grey line shows our growth function applied to the normal temperature profile without the $T_{\mathrm{eff}}$ correction. Under normal conditions our model closely tracks the TGC reference ($-3\%$ at month~12). The warm scenario produces the largest divergence: the TGC model, which does not account for reduced growth above the thermal optimum, reaches 205 tonnes while our model levels off at 148 tonnes. The effect of the $T_{\mathrm{eff}}$ correction is visible in the gap between the dotted grey and solid green lines from August onward, when sea temperatures exceed $T_{\mathrm{opt}}$.}
    \label{fig:biomass_trajectory}
\end{figure}
 
\newpage
%------------------------------------------------------------
\subsection{Model Description}
%------------------------------------------------------------
 
\subsubsection{Objective}
 
\noindent \textbf{Decision Variables} 
\\
The model has two classes of decision variables. The first is the stocking quantity $q_{u,s}^{\xi}$, representing the number of smolt stocked for cohort $(u,s)$ under scenario $\xi$. As smolt are purchased in large quantities (tens of thousands per cohort), $q_{u,s}^{\xi}$ is treated as a continuous variable. The second class consists of two sets of binary variables: $z_{u,s}^{\xi}$ indicating whether a stocking event occurs for cohort $(u,s)$, and $h_{u,s,t}^{\xi}$ indicating whether a harvest event occurs for cohort $(u,s)$ in period $t$.
 
\noindent \textbf{Objective Function} 
\\
The objective for the stochastic production planning model is to maximize the expected Net Present Value (NPV) over the five-year planning horizon. The discount factor for month $t$ is $d_t = (1 + r_m)^{-t}$ for $t \in \{0, \ldots, T-1\}$, where $r_m = (1.10)^{1/12} - 1$ is the monthly equivalent of a 10\% annual discount rate. We extend this definition to $d_T = (1 + r_m)^{-T}$ for use in the terminal-value term, which prices any biomass still standing at the end of the horizon ($t = T-1$) one period into the future to mitigate end-of-horizon effects.
 
The objective function $Z$ represents the probability-weighted sum of scenario-specific cash flows across the scenario set $\Xi$. For each scenario $\xi \in \Xi$ with probability $\pi_\xi$, where $\sum_\xi \pi_\xi = 1$, the net cash flow is derived from four primary components, all discounted by $d_t$. The first term represents harvest revenue, calculated from the harvested biomass $H^{\xi}_{\mathrm{biom}}$ and the size-dependent price $p^\xi$. The second term captures the terminal value of standing biomass $\tilde{H}^{\xi}_{\mathrm{biom}}$ at the end of the horizon, valued at a fixed terminal price $p^{\mathrm{TV}}$. The third and fourth terms represent the primary operational expenditures: the cost of stocked smolt $c^{\mathrm{smolt}}$ and the cumulative feed costs $c^{\mathrm{feed}}$ required to maintain the standing biomass $B^{\xi}$, respectively.
 
The objective function (described properly with indices and set affiliation in Section~\ref{sec:compact}) takes the following form:
 
\begin{equation}
  \max Z = \sum_{\xi \in \Xi} \pi_\xi \Big[ 
    \underbrace{\textstyle\sum d_t \  p^\xi\ H^{\xi}_{\mathrm{biom}}}_{\text{harvest revenue}}
    + \underbrace{d_T\ p^{\mathrm{TV}} \textstyle\sum \tilde{H}^{\xi}_{\mathrm{biom}}}_{\text{terminal value}}
    - \underbrace{\textstyle\sum d_t c^{\mathrm{smolt}} q^{\xi}}_{\text{smolt cost}}
    - \underbrace{\textstyle\sum d_t c^{\mathrm{feed}} B^{\xi}}_{\text{feed cost}} \Big]
  \label{eq:obj_sketch}
\end{equation}
 
 
\noindent \textbf{Pricing} 
\\
We propose the following adjustment of the historical average price of each weight class to account for the risk factors of price volatility and low liquidity discussed in Section~\ref{sec:salmonfarming} (Figure~\ref{fig:price_vol_std}). Let $\bar{p}_c$ denote the historical average price of weight class~$c$, $\sigma_c$ its historical standard deviation, and $v_c$ its share of total traded volume. The effective price used in the model is
\begin{equation}
\label{eq:adjusted_price}
  p_c^{\mathrm{eff}} = \bar{p}_c \cdot v_c^{\,\gamma} - \lambda\,\sigma_c + k,
\end{equation}
where $\gamma > 0$ controls how strongly the volume share dampens the price of thinly traded classes, and $\lambda \geq 0$ penalizes price volatility. The constant $k$ is used to calibrate the adjusted prices to realistic levels. The exact values of $k$, $\gamma$ and $\lambda$ are adjustable tuning parameters; in the base case we set $k = 5$, $\gamma = 0.02$ and $\lambda = 2$.
 
\begin{table}[ht]
\centering
\small
\begin{tabular}{l*{9}{c}}
\toprule
 & 1--2\,kg & 2--3 & 3--4 & 4--5 & 5--6 & 6--7 & 7--8 & 8--9 & 9+ \\
\midrule
Historical avg.\ (NOK/kg) & 46.08 & 54.43 & 60.59 & 62.69 & 64.92 & 67.68 & 68.49 & 69.26 & 69.01 \\
Risk-adjusted (NOK/kg) & 39.72 & 52.66 & 60.73 & 63.33 & 64.55 & 64.14 & 62.85 & 61.32 & 58.80 \\
\bottomrule
\end{tabular}
\caption{Historical average and risk-adjusted prices by weight class. The adjustment narrows the spread and penalises classes with low traded volume or high price volatility.}
\label{tab:price_adjustment}
\end{table}
 
Mid-range fish (4--6\,kg), which dominate trading volume and exhibit low price variance, retain prices close to their historical mean, while extreme weight classes are adjusted downward. This approach reflects the practical reality that selling a 4\,kg fish is considerably easier than selling a 9\,kg fish, and it ensures the model accounts for the liquidity constraints and price volatility that characterise extreme weight classes.
 
\noindent \textbf{Price-class mapping} \\
Since all fish in a cohort follow the same deterministic weight trajectory $\mathrm{wpf}^{\xi}_{s,t}$ under a given scenario, the harvest price for cohort $(u,s)$ harvested at period $t$ is determined by a direct lookup: the cohort receives price $p_c$ for the unique weight class $c \in \mathcal{C}$ satisfying $W_c \leq \mathrm{wpf}^{\xi}_{s,t} < W_{c+1}$. We write this as
\begin{equation}
    p^{\xi}_{s,t} = p_c
    \quad \text{where } c \text{ is the unique class such that }
    W_c \leq \mathrm{wpf}^{\xi}_{s,t} < W_{c+1}.
\end{equation}
As $\mathrm{wpf}^{\xi}_{s,t}$ is precomputed before the optimisation is solved, $p^{\xi}_{s,t}$ is a deterministic parameter.
 
 
\subsection{Constraints}
\subsubsection{Operational Constraints}
\label{sec:operational}
 
To track the presence of fish within a unit, we define an auxiliary binary variable ``alive'' $a^\xi_{u,s,t} \in \{0,1\}$. The first index, $u$, defines the unit in which the variable applies. The second, $s$, indicates the cohort by its starting month. The third, $t$, shows the month for which $a^\xi_{u,s,t}$ indicates presence. A cohort $(u,s)$ starts to be alive in month $t = s$ if it is stocked, $z^\xi_{u,s} = 1$:
\begin{align}
  a^\xi_{u,s,s} &= z^\xi_{u,s} && \forall\, u,\ s,\ \xi
\end{align}
 
In subsequent periods, a cohort remains alive ($a^\xi_{u,s,t} = 1$) unless a harvest event occurs in the preceding month. Once a cohort is harvested, it can no longer remain alive in the cage. If the cohort $(u,s)$ is harvested ($h^\xi_{u,s,t} = 1$) at any point $t$, then the cohort $(u,s)$ is no longer alive after that:
\begin{align}
  a^\xi_{u,s,t} &= a^\xi_{u,s,t-1} - h^\xi_{u,s,t-1} && \forall\,(u,s,t)\in \mathcal{A},\ t>s,\ \xi
\end{align}
 
To ensure that cohorts are not mixed within a single unit, we enforce a capacity constraint allowing at most one alive cohort per cage at any given time:
\begin{align}
  \sum_{s:\,(u,s,t)\in \mathcal{A}} a^\xi_{u,s,t} &\leq 1 \label{eq:C4b} && \forall\, u, \ t,\ \xi
\end{align}
The set $\mathcal{A} = \{(u,s,t) : s \leq t \leq s + T^{\mathrm{mat}},\ u \in U,\ s,t \in \mathcal{T}\}$ collects all triples $(u,s,t)$ for which a cohort stocked in unit~$u$ at period~$s$ could conceivably still be in the water at period~$t$, given the sexual maturation cap $T^{\mathrm{mat}}$.
 
Every harvest event must correspond to a specific stocking decision. While the model allows for a cohort to be stocked and harvested within the planning horizon, a harvest cannot occur without a preceding stocking event. This is enforced by:
\begin{align}
  \sum_{t:\,(u,s,t)\in \mathcal{H}^\xi} h^\xi_{u,s,t} \leq z^\xi_{u,s} && \forall\, u\in U,\ s,\ \xi
\end{align}
 
The set $\mathcal{H}^\xi = \{(u,s,t) \in \mathcal{A} : \mathrm{wpf}^\xi_{s,t} \geq \mathrm{wpf}^{\min}_{\mathrm{harv}}\}$ holds all harvest-eligible triples in $\mathcal{A}$ under scenario $\xi$. The subset $\mathcal{H}^\xi \subseteq \mathcal{A}$ only includes cohorts at months where the individual fish weight reaches the minimum required to be compatible with harvest infrastructure such as wellboat transport and slaughter facilities. Because $\mathrm{wpf}^\xi_{s,t}$ depends on the temperature scenario, $\mathcal{H}^\xi$ is scenario-specific, and harvest variables $h^\xi_{u,s,t}$ are therefore only defined for $(u,s,t) \in \mathcal{H}^\xi$.
 
 
\subsubsection{Stocking Constraints and Linearisation}
 
The stocking quantity $q^\xi_{u,s}$ is bounded above by the maximum capacity $Q^{\max}_{u}$ of unit $u$, and is forced to zero whenever no stocking event occurs ($z^\xi_{u,s} = 0$):
\begin{align}
  q^\xi_{u,s} &\leq Q^{\max}_{u}\, z^\xi_{u,s}
\end{align}
 
This constraint links the continuous headcount to the binary stocking decision via the unit-dependent upper bound:
\begin{align}
Q^{\max}_{u} = \frac{\rho^{\max} \cdot V_u}{\text{smolt size}}.
\end{align}
 
The standing biomass $B^\xi_{u,s,t}$ is derived from the biological modelling section, \ref{sec:biology}. It is given as the product of the biomass trajectory per stocked fish $\mathrm{sbpf}^\xi_{s,t}$, the continuous variable $q^\xi_{u,s}$, and the binary variable $a^\xi_{u,s,t}$ for all possible cohorts each month, $(u,s,t) \in \mathcal{A}$:
\begin{equation}
  B^\xi_{u,s,t} =
  \mathrm{sbpf}^\xi_{s,t} \cdot q^\xi_{u,s}\cdot a^\xi_{u,s,t}
  \label{eq:biom} 
\end{equation}
 
Since harvest is modelled as a single all-or-nothing event, having $h^\xi_{u,s,t} = 1$ makes $a^\xi_{u,s,t+1} = 0$ and the entire biomass in month $t$ of cohort $(u,s)$ is removed at once. Because $\mathrm{sbpf}^\xi_{s,t}$ already incorporates cumulative survival from stocking at $s$ to period $t$, using the stocked amount gives us the true biomass value. The harvested biomass $H^\xi_{\mathrm{biom},u,s,t}$ follows the same logic, but with $h^\xi_{u,s,t}$ instead of $a^\xi_{u,s,t}$. It belongs to the set of possible harvesting months $\mathcal{H}^\xi$, and only has positive values during a harvesting month ($h^\xi_{u,s,t} = 1$):
\begin{equation}
    H^\xi_{\mathrm{biom},u,s,t} = \mathrm{sbpf}^\xi_{s,t} \cdot q^\xi_{u,s} \cdot h^\xi_{u,s,t}.
\end{equation}
 
While this is the most natural representation of biomass, it is not directly implementable as the products of binary and continuous variables yield bilinear, and consequently nonlinear terms.
 
\noindent \textbf{Linearisation} \\
Both $B^\xi_{u,s,t}$ and $H^\xi_{\mathrm{biom},u,s,t}$ contain products of a binary variable ($a^\xi_{u,s,t}$ or $h^\xi_{u,s,t}$) and a continuous variable ($q^\xi_{u,s}$), and must be linearised for the MILP. We apply big-$M$ linearisation. This is enforced by:
\begin{align}
  B^\xi_{u,s,t} &\leq M_{u} \cdot a^\xi_{u,s,t}, \label{eq:lin1}\\
  B^\xi_{u,s,t} &\leq \mathrm{sbpf}^\xi_{s,t} \cdot q^\xi_{u,s}, \label{eq:lin2}\\
  B^\xi_{u,s,t} &\geq \mathrm{sbpf}^\xi_{s,t} \cdot q^\xi_{u,s} - M_{u}(1 - a^\xi_{u,s,t}), \label{eq:lin3}
\end{align}
where $M_{u}$ is a per-cohort upper bound on biomass, derived from the cage density constraint: $M_{u} = \rho^{\max} \cdot V_u$. When $a^\xi_{u,s,t} = 1$, constraints~\eqref{eq:lin2} and~\eqref{eq:lin3} force $B^\xi_{u,s,t} = \mathrm{sbpf}^\xi_{s,t} \cdot q^\xi_{u,s}$; when $a^\xi_{u,s,t} = 0$, the non-negativity of $B^\xi_{u,s,t}$ together with constraint~\eqref{eq:lin1} forces $B^\xi_{u,s,t} = 0$.
 
The same linearisation applies to $H^\xi_{\mathrm{biom},u,s,t}$ with $h^\xi_{u,s,t}$ replacing $a^\xi_{u,s,t}$:
\begin{align}
  H^\xi_{\mathrm{biom},u,s,t} &\leq \mathrm{sbpf}^\xi_{s,t} \cdot q^\xi_{u,s} \\
  H^\xi_{\mathrm{biom},u,s,t} &\leq M_u\, h^\xi_{u,s,t} \\
  H^\xi_{\mathrm{biom},u,s,t} &\geq \mathrm{sbpf}^\xi_{s,t} \cdot q^\xi_{u,s} - M_u\,(1 - h^\xi_{u,s,t})
\end{align}
 
\begin{figure}[ht]
    \centering
    \includegraphics[width=1.0\linewidth]{bilder/cohort.png}
    \caption{Representation of the lifespan of a single cohort. This example shows a cohort stocked in month 3 and harvested in month 20, followed by a two-month fallow period for the cohort's cage. The alive indicator $a^\xi_{u,s,t}$ is zero before stocking and after harvest, and equals one inside this interval.}
    \label{fig:cohort_lifespan}
\end{figure}
 
\subsubsection{Regulatory Constraints}
 
\noindent \textbf{Fallowing} \\
Norwegian aquaculture regulation (\emph{Akvakulturdriftsforskriften} §40) mandates a minimum two-month fallow period after each production cycle. This implies that for a stocking event $z^\xi_{u,t} = 1$ to occur, cage $u$ must have been empty for the two preceding periods, with $h^\xi_{u,s,t-1} = 0$ and $h^\xi_{u,s,t-2} = 0$. Conversely, any harvest event requires the cage to remain fallow for the subsequent two months.
 
To enforce this, we sum over all possible cohorts $(u,s)$ that could be present in the unit. This ensures that if any cohort was harvested in the previous two months, a new stocking is prohibited:
 
\vspace{1em}
 
\noindent\emph{Unit-level fallowing} (two-month post-harvest):
\begin{equation}
  z^\xi_{u,t} + \sum_{s:\,(u,s,t)\in \mathcal{H}^\xi} h^\xi_{u,s,t-1} + \sum_{s:\,(u,s,t)\in \mathcal{H}^\xi} h^\xi_{u,s,t-2} \leq 1 \qquad \forall\, u,t,\xi \label{eq:C5}
\end{equation}
 
\vspace{1em}
 
\noindent In addition to unit-level requirements, some areas are subject to coordinated lice management zones (\emph{lusesoner}). As prescribed by \citet{mattilsynet_brakklegging2} and \citet{vetinst_brakklegging}, these regulations require all facilities within 5\,km of each other to fallow simultaneously for a fixed period, typically four weeks, which we round of to one month.
 
We enforce this by requiring that no alive cohorts exist at any unit during its location's coordinated fallow period $\mathcal{F}_{\ell(u)}$:
\begin{align}
  a^\xi_{u,s,t} &= 0
    && \forall\, (u,s,t) \in \mathcal{A} \text{ with } t \in \mathcal{F}_{\ell(u)},\ \xi
    \label{eq:C6a}
\end{align}
where $\mathcal{F}_{\ell}$ is the set of scheduled coordinated fallow months for location $\ell \in \mathcal{L}$.
 
\begin{figure}[ht]
    \centering
    \includegraphics[width=1\linewidth]{bilder/fallowing_violations.png}
    \caption{Subsections of a planning period for four units under two different schedules. (a) shows a compliant plan where every cohort respects both fallowing rules: \textit{unit-level fallowing} and the \textit{location's coordinated fallow period}. Plan (b) violates the first rule for unit one and the second rule for unit three.}
    \label{fig:fallowing_violations}
\end{figure}
 
\noindent \textbf{Maximal Allowed Biomass} \\
Norwegian salmon farming is governed by a hierarchical system of biomass constraints that operate simultaneously at three spatial levels. At the cage level, the density constraint limits biomass to $\rho^{\max} \cdot V_u$, where $\rho^{\max} = 25\;\mathrm{kg/m}^3$ and $V_u$ is the volume of cage $u$. At the location level, each facility holds a licensed MAB limit $\mathrm{MAB}_\ell$, which caps the total standing biomass across all cages at that location. At the regional level, an aggregate cap $\mathrm{MAB}_{\mathrm{reg}}$ governs the total biomass within the license zone, currently set at 35,000 tonnes in the parameterization used here.
 
\begin{figure}[ht]
    \centering
    \includegraphics[width=1\linewidth]{bilder/production_hierarchy.png}
    \caption{The hierarchy of regulatory production limits. On the lowest level we have density limits per cage $u$, followed by per-location biomass limits with several cages, and at the highest order the biomass limit for a region with several locations.}
    \label{fig:production_hierarchy}
\end{figure}
 
The natural formulation for these inequalities are:
\begin{align}
 \sum_{s:\,(u,s,t)\in\mathcal{A}} B^\xi_{u,s,t} &\leq \rho^{\max} V_u 
   && \forall\, u \in U,\ t \in \mathcal{T},\ \xi \in \Xi \\
  \sum_{u:\,\ell(u)=\ell} \sum_{s:\,(u,s,t)\in\mathcal{A}} B^\xi_{u,s,t} 
   &\leq \mathrm{MAB}_\ell 
   && \forall\, \ell \in \mathcal{L},\ t \in \mathcal{T},\ \xi \in \Xi \\
  \sum_{u\in U} \sum_{s:\,(u,s,t)\in\mathcal{A}} B^\xi_{u,s,t} 
   &\leq \mathrm{MAB}_{\mathrm{reg}}
   && \forall\, t \in \mathcal{T},\ \xi \in \Xi
\end{align}
We instead implement these constraints in elastic form, allowing the limits to be exceeded at a per-kilogram penalty $c^{\mathrm{pen}}$. For each hard constraint we introduce a non-negative violation variable that measures the amount by which the limit is breached:
\begin{align}
  \sum_{s:\,(u,s,t)\in\mathcal{A}} B^\xi_{u,s,t}
    &\leq \rho^{\max} V_u + v^{\mathrm{U},\xi}_{u,t} \\
  \sum_{u:\,\ell(u)=\ell}\ \sum_{s:\,(u,s,t)\in\mathcal{A}} B^\xi_{u,s,t}
    &\leq \mathrm{MAB}_\ell + v^{\mathrm{L},\xi}_{\ell,t} \\
  \sum_{u \in U}\ \sum_{s:\,(u,s,t)\in\mathcal{A}} B^\xi_{u,s,t}
    &\leq \mathrm{MAB}_{\mathrm{reg}} + v^{\mathrm{R},\xi}_{t}
\end{align}
with $v^{\mathrm{U},\xi}_{u,t},\, v^{\mathrm{L},\xi}_{\ell,t},\, v^{\mathrm{R},\xi}_{t} \in \mathbb{R}_{\geq 0}$, and add a penalty term to the objective:
\begin{equation}
  -\, c^{\mathrm{pen}} \sum_{\xi \in \Xi} \pi_\xi
    \left[
      \sum_{u \in U} \sum_{t \in \mathcal{T}} v^{\mathrm{U},\xi}_{u,t}
      \;+\; \sum_{\ell \in \mathcal{L}} \sum_{t \in \mathcal{T}} v^{\mathrm{L},\xi}_{\ell,t}
      \;+\; \sum_{t \in \mathcal{T}} v^{\mathrm{R},\xi}_{t}
    \right].
  \label{eq:pen}
\end{equation}
The motivation for this elastic formulation is given in Section~\ref{sec:stochastic_programming}.
 
\begin{figure}[ht]
    \centering
    \includegraphics[width=1\linewidth]{bilder/mab_violation.png}
    \caption{Aggregate regional biomass $\sum_u\sum_s B^{\xi}_{u,s,t}$ over the planning horizon under two stocking strategies. (a) A staggered plan keeps biomass comfortably below the regional MAB ceiling throughout. (b) A synchronised plan produces two problematic peaks: the first breaches MAB inside the planning horizon, while the second stays feasible across $[0, T{-}1]$ but would in reality continue growing once the horizon ends and breach MAB shortly thereafter. Without a constraint on biomass carried to the terminal period, the model has no incentive to prevent this kind of post-horizon violation, the motivation for the terminal-biomass cap~\eqref{eq:terminal_biomass_cap}.}
    \label{fig:mab_violation}
\end{figure}
 
\subsubsection{End of horizon constraints}
\noindent \textbf{End-of-horizon treatment} \\
To complement the terminal harvest logic, we impose an upper bound on terminal standing biomass to avoid a plan that would breach density and biomass constraints right after the horizon window as the premature fish grow into harvesting size. The sum of biomass in all cohorts at the end of the horizon ($t = T - 1$) for each scenario $\xi$ should be less than or equal to the maximum of either the initial biomass at horizon start or a fraction $k$ of the regional biomass limit (e.g.\ 90\%):
\begin{equation}
\sum_{u \in U}\;\sum_{s:\,(u,s,T{-}1)\in\mathcal{A}} B^\xi_{u,s,T-1}
    \leq \max \Big\{\sum_{u \in U^0}\;B^\xi_{u,0,0}, \ \  k \cdot \mathrm{MAB}_{\mathrm{reg}}\Big\}
    \qquad \forall\, \xi \in \Xi,
  \label{eq:terminal_biomass_cap}
\end{equation}
where $U^{0} \subseteq U$ is the set of units with fish at $t=0$, and $k \in (0,1)$ is a tuning parameter.
 
The terminal harvest mechanism and the terminal biomass cap together prevent two problems: no new stocking near the horizon end (addressed by the terminal price) and the excessive accumulation of biomass at the horizon end (addressed by constraint~\eqref{eq:terminal_biomass_cap}).
 
 
%------------------------------------------------------------
\subsection{Uncertainty representation and non-anticipativity constraints}
\label{sec:uncertainty}
%------------------------------------------------------------
 
\subsubsection{Scenario tree}
The planning horizon of 60 months is divided into four stages of 15 months each:
\begin{equation*}
  \mathcal{T}_0 = [0,14],\quad
  \mathcal{T}_1 = [15,29],\quad
  \mathcal{T}_2 = [30,44],\quad
  \mathcal{T}_3 = [45,59].
\end{equation*}
At each stage, one of three temperature realisations $\{\text{cold, normal, warm}\}$ is drawn. Growth depends proportionally on the temperature of the given stage. Mortality is path-dependent: two consecutive warm stages trigger an elevated monthly mortality rate during the second. Stages are assumed independent, so each contributes three possible outcomes, yielding $|\Xi| = 3^4 = 81$ equally probable leaf scenarios with probability $1/81$.
 
\begin{figure}[ht]
    \centering
    \includegraphics[width=1\linewidth]{bilder/tree.png}
    \caption{Scenario tree illustrating the four-stage temperature realisation process, which drives growth and mortality across all 81 leaf scenarios.}
    \label{fig:scenario_tree}
\end{figure}
 
 
\subsubsection{Non-anticipativity}
As we are dealing with several future state spaces, we need to make sure that the model does not anticipate a supposedly unknown future. All decision variables should only take into consideration the history up to that point in time, as well as the probability-weighted distribution of future outcomes. However, the scenario tree structure lets the model know which scenario branch it is on, and therefore also what future events will occur. To prevent this foresight from contaminating decisions, we introduce non-anticipativity constraints. These force decision variables associated with stage k to be identical across any two branches that share the same history up to stage $k$, even though they diverge afterward.
 
For any two scenarios $\xi, \xi' \in \Xi$ sharing the same node at stage $k$ (written $\xi \sim_k \xi'$, defined as agreeing on temperature realisations for stages $0, 1, \ldots, k$), all decisions taken during stage~$k$ must agree:
\begin{align}
  z^\xi_{u,s}   &= z^{\xi'}_{u,s}
    && \forall\, u \in U,\ s \in \mathcal{T}_k,\ k \in \mathcal{K},\ \xi\sim_k\xi' \label{eq:NA1}\\
  q^\xi_{u,s}   &= q^{\xi'}_{u,s}
    && \forall\, u \in U,\ s \in \mathcal{T}_k,\ k \in \mathcal{K},\ \xi\sim_k\xi' \label{eq:NA2}\\
  h^\xi_{u,s,t} &= h^{\xi'}_{u,s,t}
    && \forall\, u \in U,\ s \in \mathcal{T},\ t \in \mathcal{T}_k,\
       (u,s,t) \in \mathcal{H}^\xi \cap \mathcal{H}^{\xi'},\
       k \in \mathcal{K},\ \xi\sim_k\xi' \label{eq:NA3}
\end{align}
In~\eqref{eq:NA1}--\eqref{eq:NA2} the stocking decision is taken at period~$s$, so it is the stage containing~$s$ that determines the NA grouping. In~\eqref{eq:NA3} the harvest decision is taken at period~$t$, so the grouping depends on the stage containing~$t$ regardless of when the cohort was stocked. The auxiliary variables $a^\xi$, $B^\xi$, and $H^\xi_{\mathrm{biom}}$ inherit non-anticipativity automatically, since they are fully determined by the decision variables through the linearisation and operational constraints (Section~\ref{sec:operational}).
 
 
\newpage
%------------------------------------------------------------
\subsection{Mathematical Model}
\label{sec:compact}
%------------------------------------------------------------
 
We present the mathematical formulation of the salmon farming SP in the extensive form. The model couples one copy of the single-scenario MILP per scenario $\xi \in \Xi$ through non-anticipativity constraints. 
 
 
\subsubsection{Sets and indices}
\begin{align*}
  U        &: \text{production units (cages)},\quad U^{0} \subseteq U\ \text{units with fish at } t=0 \\
  \mathcal{L} &: \text{set of locations},\quad \ell\colon U\to\mathcal{L},\ \ell(u)\ \text{denotes the location of unit } u \\
  \mathcal{T} &= \{0,\dots,T{-}1\},\ \text{monthly planning periods} \\
  \mathcal{K} &= \{0,1,2,3\},\ \text{decision stages, with month ranges} \\
              &\quad \mathcal{T}_0 = [0,14],\ \mathcal{T}_1 = [15,29],\ \mathcal{T}_2 = [30,44],\ \mathcal{T}_3 = [45,59] \\
  \mathcal{A} &= \{(u,s,t) : u \in U,\ s,t \in \mathcal{T},\ s \leq t \leq s + T^{\mathrm{mat}}\} \\
  \mathcal{H}^\xi    &= \{(u,s,t) \in \mathcal{A} : \mathrm{wpf}^\xi_{s,t} \geq \mathrm{wpf}^{\min}_{\mathrm{harv}}\},\quad \xi \in \Xi \\
  \Xi      &: \text{scenario set with probabilities } \pi_\xi > 0,\quad \textstyle\sum_{\xi \in \Xi} \pi_\xi = 1
\end{align*}
 
We further partition the harvest set into regular and terminal harvests:
\begin{align*}
  \mathcal{H}^\xi_{\mathrm{reg}}  &= \{(u,s,t) \in \mathcal{H}^\xi : t < T{-}1 \}, \\
  \mathcal{H}^\xi_{\mathrm{term}} &= \{(u,s,t) \in \mathcal{H}^\xi : t = T{-}1 \}.
\end{align*}
 
 
\subsubsection{Parameters}
\begin{align*}
  V_u                                    &: \text{cage volume (m}^3) \\
  \rho^{\max}                            &: \text{density limit (kg/m}^3) \\
  \mathrm{MAB}_\ell,\ \mathrm{MAB}_{\mathrm{reg}} &: \text{location and regional biomass limits (kg)} \\
  Q^{\max}_{u}                           &: \parbox[t]{0.72\linewidth}{maximum stocking headcount for unit $u$, derived from the cage density limit and smolt size as $Q^{\max}_u = \rho^{\max} V_u / \text{smolt size}$} \\
  c^{\mathrm{smolt}}                     &: \text{smolt cost (NOK/head)} \\
  c^{\mathrm{feed}}                      &: \text{feed cost (NOK/kg standing biomass/month)} \\
  p^\xi_{s,t}                            &: \parbox[t]{0.72\linewidth}{harvest price (NOK/kg), precomputed per scenario from the weight-class rule} \\
  p^{\mathrm{TV}}                        &: \text{terminal value (NOK/kg)} \\
  d_t                                    &: \text{monthly discount factor at period } t \\
  \mathrm{sbpf}^\xi_{s,t}                &: \parbox[t]{0.72\linewidth}{survival-adjusted biomass per fish (kg) for a cohort stocked at $s$, evaluated at $t$ under scenario $\xi$} \\
  \mathrm{wpf}^\xi_{s,t}                 &: \text{individual fish weight (g) at $t$ for cohort stocked at $s$ under scenario } \xi \\
  \mathrm{wpf}^{\min}_{\mathrm{harv}}    &: \text{minimum fish weight for harvest eligibility (g)} \\
  T^{\mathrm{mat}}                       &: \text{maturation cap (months)} \\
  M_u                                    &: \text{big-}M \text{ bound per unit},\quad M_u = \rho^{\max} V_u \\
  \mathcal{F}_\ell                       &: \text{scheduled fallow months for location } \ell \\
  c^{\mathrm{pen}}                       &: \text{penalty for biomass-constraint violation (NOK/kg)}
\end{align*}
 
\subsubsection{Decision variables}
 
Each scenario $\xi \in \Xi$ carries its own copy of all variables:
 
\paragraph{Stocking and inventory.}
\begin{align*}
  z^\xi_{u,s} &\in \{0,1\} && \forall\, u \in U,\ s \in \mathcal{T},\ \xi \in \Xi \\
    &\quad \text{1 if a cohort is stocked in unit } u \text{ at period } s \\
  q^\xi_{u,s} &\in \mathbb{R}_{\geq 0} && \forall\, u \in U,\ s \in \mathcal{T},\ \xi \in \Xi \\
    &\quad \text{number of smolts stocked}
\end{align*}
 
\paragraph{Harvest and cohort dynamics.}
\begin{align*}
  h^\xi_{u,s,t} &\in \{0,1\} && \forall\, (u,s,t) \in \mathcal{H}^\xi,\ \xi \in \Xi \\
    &\quad \text{1 if cohort } (u,s) \text{ is harvested at } t \\
  a^\xi_{u,s,t} &\in \{0,1\} && \forall\, (u,s,t) \in \mathcal{A},\ \xi \in \Xi \\
    &\quad \text{1 if cohort } (u,s) \text{ is alive at } t \\
  B^\xi_{u,s,t} &\in \mathbb{R}_{\geq 0} && \forall\, (u,s,t) \in \mathcal{A},\ \xi \in \Xi \\
    &\quad \text{total biomass of cohort } (u,s) \text{ at } t \\
  H^\xi_{\mathrm{biom},u,s,t} &\in \mathbb{R}_{\geq 0} && \forall\, (u,s,t) \in \mathcal{H}^\xi,\ \xi \in \Xi \\
    &\quad \text{total harvested biomass of cohort } (u,s) \text{ at } t
\end{align*}
 
\paragraph{Violation variables.}
\begin{align*}
  v^{\mathrm{U},\xi}_{u,t} &\in \mathbb{R}_{\geq 0} && \forall\, u \in U, t \in \mathcal{T}, \xi \in \Xi \\
    &\quad \text{unit-density constraint violation at } u, t \\
  v^{\mathrm{L},\xi}_{\ell,t} &\in \mathbb{R}_{\geq 0} && \forall\, \ell \in \mathcal{L}, t \in \mathcal{T}, \xi \in \Xi \\
    &\quad \text{location MAB constraint violation at } \ell, t \\
  v^{\mathrm{R},\xi}_{t} &\in \mathbb{R}_{\geq 0} && \forall\, t \in \mathcal{T}, \xi \in \Xi \\
    &\quad \text{regional MAB constraint violation at } t
\end{align*}
 
\subsubsection{Objective}
Maximise expected net present value:
\begin{align}
  \max\ Z = \sum_{\xi \in \Xi} \pi_\xi \Bigg[
    &\phantom{+}\sum_{(u,s,t)\in \mathcal{H}^\xi_{\mathrm{reg}}} d_t\, p^\xi_{s,t}\, H^{\xi}_{\mathrm{biom},u,s,t} \notag\\
    &+\, d_T\, p^{\mathrm{TV}} \sum_{(u,s,t)\in \mathcal{H}^\xi_{\mathrm{term}}} H^{\xi}_{\mathrm{biom},u,s,t} \notag\\
    &-\sum_{u \in U}\sum_{s \in \mathcal{T}} d_s\, c^{\mathrm{smolt}}\, q^\xi_{u,s} \notag\\
    &-\sum_{u \in U}\sum_{t \in \mathcal{T}} d_t\, c^{\mathrm{feed}}
      \sum_{s:\,(u,s,t)\in\mathcal{A}} B^\xi_{u,s,t} \notag \\
    &-\, c^{\mathrm{pen}}\bigg(\sum_{u \in U} \sum_{t \in \mathcal{T}} v^{\mathrm{U},\xi}_{u,t} + \sum_{\ell \in \mathcal{L}} \sum_{t \in \mathcal{T}} v^{\mathrm{L},\xi}_{\ell,t} + \sum_{t \in \mathcal{T}} v^{\mathrm{R},\xi}_{t} \bigg) \Bigg]
  \label{eq:obj}
\end{align}
 
\medskip
\subsubsection{Constraints}
 
\noindent \textbf{Stocking and headcount.}
\begin{align}
  q^\xi_{u,s} &\leq Q^{\max}_{u}\, z^\xi_{u,s}
    && \forall\, u \in U,\ s \in \mathcal{T},\ \xi \in \Xi \label{eq:C1} \\
  \sum_{t:\,(u,s,t)\in \mathcal{H}^\xi} h^\xi_{u,s,t} &\le z^\xi_{u,s}
    && \forall\, u \in U,\ s \in \mathcal{T},\ \xi \in \Xi \label{eq:C2}
\end{align}
 
 
\noindent \textbf{Alive-status propagation.}
\begin{align}
  a^\xi_{u,s,s} &= z^\xi_{u,s}
    && \forall\, u \in U,\ s \in \mathcal{T},\ \xi \in \Xi \label{eq:C3a}\\
  a^\xi_{u,s,t} &= a^\xi_{u,s,t-1} - h^\xi_{u,s,t-1}
    && \forall\, (u,s,t) \in \mathcal{A},\ t > s,\ \xi \in \Xi \label{eq:C3b}
\end{align}
where $h^\xi_{u,s,t-1} = 0$ whenever $(u,s,t{-}1) \notin \mathcal{H}^\xi$, since the corresponding variable is not defined.
 
\noindent \textbf{One cohort per unit.}
\begin{align}
  \sum_{s:\,(u,s,t)\in \mathcal{A}} a^\xi_{u,s,t} &\leq 1
    && \forall\, u \in U,\ t \in \mathcal{T},\ \xi \in \Xi \label{eq:C4}
\end{align}
 
 
\noindent \textbf{Fallow constraints.}
 
\noindent\emph{Unit-level} (post-harvest fallow):
\begin{equation}
  z^\xi_{u,t}
  + \sum_{s:\,(u,s,t{-}1)\in \mathcal{H}^\xi} h^\xi_{u,s,t-1}
  + \sum_{s:\,(u,s,t{-}2)\in \mathcal{H}^\xi} h^\xi_{u,s,t-2}
  \leq 1
  \qquad \forall\, u \in U,\ t \in \mathcal{T},\ \xi \in \Xi \label{eq:C5_compact}
\end{equation}
 
 \noindent\emph{Location-level} (scheduled):
\begin{align}
  a^\xi_{u,s,t} &= 0
    && \forall\, (u,s,t) \in \mathcal{A},\ \text{with } t \in \mathcal{F}_{\ell(u)},\ \xi \in \Xi \label{eq:C6}
\end{align}
 
 
\noindent \textbf{Biomass linearisation.} \\
The standing biomass of a cohort satisfies $B^\xi_{u,s,t} = \mathrm{sbpf}^\xi_{s,t} \cdot q^\xi_{u,s} \cdot a^\xi_{u,s,t}$, which we express through the following linear constraints:
\begin{align}
  B^\xi_{u,s,t} &\leq \mathrm{sbpf}^\xi_{s,t} \cdot q^\xi_{u,s}
    && \forall\, (u,s,t) \in \mathcal{A},\ \xi \in \Xi \\
  B^\xi_{u,s,t} &\leq M_u\, a^\xi_{u,s,t}
    && \forall\, (u,s,t) \in \mathcal{A},\ \xi \in \Xi \\
  B^\xi_{u,s,t} &\geq \mathrm{sbpf}^\xi_{s,t} \cdot q^\xi_{u,s}
    - M_u\,(1 - a^\xi_{u,s,t})
    && \forall\, (u,s,t) \in \mathcal{A},\ \xi \in \Xi 
\end{align}
Analogously, the harvested biomass satisfies $H^\xi_{\mathrm{biom},u,s,t} = \mathrm{sbpf}^\xi_{s,t} \cdot q^\xi_{u,s} \cdot h^\xi_{u,s,t}$, expressed as:
\begin{align}
  H^\xi_{\mathrm{biom},u,s,t} &\leq \mathrm{sbpf}^\xi_{s,t} \cdot q^\xi_{u,s}
    && \forall\, (u,s,t) \in \mathcal{H}^\xi,\ \xi \in \Xi \\
  H^\xi_{\mathrm{biom},u,s,t} &\leq M_u\, h^\xi_{u,s,t}
    && \forall\, (u,s,t) \in \mathcal{H}^\xi,\ \xi \in \Xi \\
  H^\xi_{\mathrm{biom},u,s,t} &\geq \mathrm{sbpf}^\xi_{s,t} \cdot q^\xi_{u,s}
    - M_u\,(1 - h^\xi_{u,s,t})
    && \forall\, (u,s,t) \in \mathcal{H}^\xi,\ \xi \in \Xi 
\end{align}
 
 
\noindent \textbf{Capacity and regulatory constraints.}
\begin{align}
  \sum_{s:\,(u,s,t)\in\mathcal{A}} B^\xi_{u,s,t}
    &\leq \rho^{\max} V_u + v^{\mathrm{U},\xi}_{u,t}
    && \forall\, u \in U,\ t \in \mathcal{T},\ \xi \in \Xi \label{eq:C7s}\\
  \sum_{u:\,\ell(u)=\ell}\ \sum_{s:\,(u,s,t)\in\mathcal{A}} B^\xi_{u,s,t}
    &\leq \mathrm{MAB}_\ell + v^{\mathrm{L},\xi}_{\ell,t}
    && \forall\, \ell \in \mathcal{L},\ t \in \mathcal{T},\ \xi \in \Xi \label{eq:C8s}\\
  \sum_{u \in U}\ \sum_{s:\,(u,s,t)\in\mathcal{A}} B^\xi_{u,s,t}
    &\leq \mathrm{MAB}_{\mathrm{reg}} + v^{\mathrm{R},\xi}_{t}
    && \forall\, t \in \mathcal{T},\ \xi \in \Xi \label{eq:C9s}
\end{align}
 
 
\noindent \textbf{Terminal biomass constraint.}
\begin{equation}
  \sum_{u \in U}\;\sum_{s:\,(u,s,T{-}1)\in\mathcal{A}} B^\xi_{u,s,T-1}
    \leq \max \Big\{\sum_{u \in U^0}\;B^\xi_{u,0,0}, \ \  k \cdot \mathrm{MAB}_{\mathrm{reg}}\Big\}
  \qquad \forall\, \xi \in \Xi  \label{eq:C10}
\end{equation}
where $k \in (0,1)$ is a tuning parameter.
 
\noindent \textbf{Non-anticipativity constraints.}
For any two scenarios $\xi, \xi'$ sharing the same node at stage $k$ (written $\xi \sim_k \xi'$), all decisions taken during stage~$k$ must agree:
\begin{align}
  z^\xi_{u,s}   &= z^{\xi'}_{u,s}
    && \forall\, u \in U,\ s \in \mathcal{T}_k,\ k \in \mathcal{K},\ \xi\sim_k\xi'
    \label{eq:NA1}\\
  q^\xi_{u,s}   &= q^{\xi'}_{u,s}
    && \forall\, u \in U,\ s \in \mathcal{T}_k,\ k \in \mathcal{K},\ \xi\sim_k\xi'
    \label{eq:NA2}\\
  h^\xi_{u,s,t} &= h^{\xi'}_{u,s,t}
    && \forall\, u \in U,\ s \in \mathcal{T},\ t \in \mathcal{T}_k,\
       (u,s,t) \in \mathcal{H}^\xi \cap \mathcal{H}^{\xi'},\
       k \in \mathcal{K},\ \xi\sim_k\xi'
    \label{eq:NA3}
\end{align}
 
 


%-----------------------------------------------------------
\subsection{Algorithm}
\label{sec:algorithm}
%-----------------------------------------------------------
The model in Section~\ref{sec:compact} is a four-stage stochastic mixed-integer linear
program (SMILP) with $3^4 = 81$ equally probable leaf scenarios. Solving its deterministic equivalent (DE) requires making a plan that takes into account 81 possible outcomes. It should have a planned response for the remainder of the horizon for each unique temperature realisation. The combination of integrality, scenario count, and a five-year horizon makes the DE computationally demanding, as we confirm empirically in Section~\ref{sec:Determinisitic_compared_Stochastic}. This section describes the matheuristic developed to solve the problem at scale.


The algorithm is a two-phase Progressive Hedging matheuristic that we refer to
throughout the remainder of the thesis as Binary-Fixing Progressive Hedging Algorithm
(B-PHA). It builds on the PH-for-MILP techniques described in
Section~\ref{sec:ph_milp}: scenario-decomposed solves with a Lagrangean penalty
driving binary variables toward consensus, variable fixing once consensus is reached,
and force-fixing variables to the value closest to consensus when convergence stagnates. Phase~1 runs the PH
loop and eventually returns a scenario-tree production calendar (stocking-and-harvest skeleton) indexed by tree node, such that any two scenarios sharing the same partial history up to
a given stage agree on all binary decisions taken at that stage. Phase~2 fixes that skeleton and
solves the resulting continuous stochastic program in extensive form, completing a good,
near-optimal solution to the full production planning problem.
Section~\ref{sec:lit_vs_own} identifies the three features specific to the salmon
production planning problem. Section~\ref{sec:binaryph} describes Phase~1 in detail.
Section~\ref{sec:lpresolve} describes Phase~2.

%-----------------------------------------------------------
\subsubsection{What is from the literature, what is our own}
\label{sec:lit_vs_own}
%-----------------------------------------------------------

Phase~1 takes inspiration from the PH-for-MILP literature described in
Section~\ref{sec:ph_milp}: consensus fixing \citep{watson2011progressive} and slamming
inspired by \citet{watson2011progressive}. The two-phase structure, in which Phase~1
commits a binary skeleton and Phase~2 solves a restricted problem with binaries fixed,
is adapted from \citet{lokketangen1996progressive} and \citet{crainic2009progressive}.
Three features of B-PHA are tailored to the salmon production planning problem rather
than inherited from the literature.


\textit{First, the binary/continuous block split.} The idea of separating integer
decisions from continuous ones and recovering the continuous variables afterwards
appears in \citet{lokketangen1996progressive} and is formalised into an explicit
two-phase structure by \citet{crainic2009progressive}. We adapt this to the salmon
production planning problem by partitioning the non-anticipative decisions specifically
into a binary block ($z$, $h$) and a continuous block ($q$), and penalising only the
binary block during Phase~1. The motivation for our partition comes from the structure
of the salmon production planning problem. The stocking quantity $q$ takes values in
the tens of thousands, and driving consensus on such a continuous variable without
convexity is hard. The binary decisions $z$ and $h$, on the other hand, live in a base-2
combinatorial space: each variable can only be $0$ or $1$, so the search space per node
is finite and small. 
%We therefore think that PH can be forced to consensus on the binary block in reasonable time, even without convergence guarantees, while the continuous is  to Phase~2. 
Once the binaries are fixed, the only
remaining free decision variable is the continuous stocking quantity $q$, and the
problem reduces to a pure stochastic LP. This is what makes Phase~2 tractable as a
single extensive-form LP across all 81 scenarios.
% We claim this as our own contribution: an adaptation of the two-phase decomposition idea from the literature, redesigned to exploit the specific structure of the salmon production planning problem.

\textit{Second, variable fixing and cycle-breaking slam.} Consensus fixing
\citep{watson2011progressive} and the general idea of slamming, meaning to force-fix
disagreeing binaries when convergence stalls, are both inherited from the literature
and adapted here. Consensus fixing fires
once every scenario at a node has agreed on the same integer value for sufficiently many
consecutive iterations, locking that variable and removing it from the augmented
objective. However, during trial runs we noticed that some binary variables cycled between values across iterations without agreeing. These oscillating variables motivated us to introduce the slamming action to avoid stagnation in solution convergence. 
%The slam targets oscillating variables by looking at the recent consensus history of each unfixed binary and fixing any variable whose consensus has oscillated by more than $\delta_{\mathrm{vol}} = 0.15$ over the last $w_{\mathrm{slam}}$ iterations, rounding its mean to the nearest integer. 
%The slam trigger monitors a rolling window of disagreement counts and max-deviation values, firing when both have stagnated while the algorithm has already made meaningful progress. 
A second way we fix deviating variables at a node is \textit{end-stagnation}: some variables neither oscillate nor converge, but simply stay stuck. A separate tail-stall escape handles
these by force-fixing them to the nearest integer once the number of disagreeing variables stays flat and small for $w_{\mathrm{tail}}$ consecutive iterations. Both the oscillating slam and the end-stagnation slam are inspired by \citet{watson2011progressive}, but are our own adaptations. 

\textit{Third, the $z$--$h$ consistency repair.} The slam mechanisms force-fix $z$
and $h$ independently, and it is precisely these force-fixed variables that can introduce
two kinds of problems into the binary consensus exiting Phase~1. The first is a
harvest variable $h^\xi_{u,s,t} = 1$ paired with a stocking variable $z_{u,s} = 0$ at
the same tree node: a cohort is scheduled to be harvested but was never stocked. The
second is the reverse: $z_{u,s} = 1$ with no corresponding harvest $h_{u,s,t} = 1$ at
any time in the plan, meaning a cohort is stocked but never harvested. The first is operationally impossible while the second is possible, but not at all optimal and could happen as a consequence of a poorly fixed variable. Hence, we add a repair pass that targets the force-fixed variables before the binary skeleton is ready. 
It is straightforward to execute
because the production plan at this stage is a simple binary array indexed by cage,
cohort, time, and tree node: checking whether a stocking decision has a matching
harvest, and vice versa, is a direct lookup with no solving required. The first pass
scans all harvest variables fixed to $1$ and zeros any whose corresponding stocking
variable is $0$; the second scans all stocking variables fixed to $1$ and zeros any for
which no harvest variable remains set to $1$ anywhere in the plan. 
%The repair works because the force-fixed variables are committed to their rounded consensus mean rather than an arbitrary value, so the resulting binary array tends to be close to a feasible production plan and contradictions are localised and easy to correct. 
There is no guarantee that the repair improves the solution, only that it makes the skeleton consistent and can pass it on to Phase~2.

The algorithm was developed through trial and error while trying to implement the PH-for-MILP algorithm on the salmon production problem. The hyperparameters specified in the upcoming section are based of our intuition and adjusted throughout development and during trial runs, when they seemed to make the algorithm run smoother and produce more profitable plans. The cost of these heuristics is that the binary skeleton committed by Phase~1 is biased by the algorithm's hyper-parameters, such as the penalty scheduling and
slamming thresholds. Generally, it is not the optimal binary solution to the underlying SMILP. We quantify this gap empirically in Section~\ref{sec:sp_vs_ef}.

%-----------------------------------------------------------
\subsubsection{Phase 1 --- Binary-Fixing Progressive Hedging}
\label{sec:binaryph}
%-----------------------------------------------------------

 \paragraph{Terminology.}
Throughout this section, ``tree node'' (or just ``node'') means a branching point in the scenario tree, identified by a partial temperature history. ``Decision variable'' means an individual non-anticipativity variable instance, e.g.~$z^\xi_{u,s}$ for a specific scenario $\xi$, cage $u$, and month $s$. Several scenarios sharing a tree node share the corresponding entries of $z, h$ at that node, and our consensus values $\bar{z}, \bar{h}$ are indexed by tree node accordingly. These are named \textit{stocking}- and \textit{harvest timing nodes} respectively. We use the term \textit{participating} for the variables if they belong to the timing node in question.
% Existing-cohort harvest timing $h^{\mathrm{exist}}$ is subject to the same non-anticipativity treatment as $h$ throughout Phase~1 and is omitted from the notation for brevity.

\paragraph{What Phase 1 produces.}
Phase~1 returns a single production calendar that is the same in every scenario sharing the same partial temperature history. Concretely, it returns binary values $\bar{z}_{u,s,n}$ and $\bar{h}_{u,s,t,n}$ indexed by tree node $n$, telling us in which months and in which cages the consensus plan stocks and harvests, conditional on the temperature outcomes observed up to that node. The continuous stocking headcount $q$ is left free for Phase~2.

\paragraph{Setup.}
Each scenario $\xi$ has its own copy of the production MILP from Section~\ref{sec:compact}, with growth rates and mortality coefficients set from $\xi$'s temperature path. We index scenarios $\xi = 1, \ldots, S$ with $S=81$ and write $\mathcal{P}^\xi$ for the per-scenario MILP. The binary non-anticipativity decisions in scenario $\xi$ are $z^\xi_{u,s}$ and $h^\xi_{u,s,t}$, each grouped by tree node so that two scenarios sharing a node $n$ share the corresponding entries. The probability-weighted node consensus values are
\begin{align}
  \bar{z}^{(k)}_{u,s,n}    &= \sum_{\xi \in \mathrm{node}(n)} p_{\xi|\mathrm{node}(n)}\, z^{\xi,(k)}_{u,s},\\
  \bar{h}^{(k)}_{u,s,t,n}  &= \sum_{\xi \in \mathrm{node}(n)} p_{\xi|\mathrm{node}(n)}\, h^{\xi,(k)}_{u,s,t},
\end{align}
where the conditional probability $p_{\xi|\mathrm{node}(n)}$ is uniform $1/|\mathrm{node}(n)|$ for this instance. The Lagrange multipliers on the two blocks are written $w^{\xi,z}_{u,s}$ and $w^{\xi,h}_{u,s,t}$. The penalty parameter is $\rho$ and the iteration counter is $k$.

\paragraph{Step 0: initial solve and penalty calibration.}
All 81 subproblems are solved independently with no penalty term, yielding initial binaries $z^{\xi,(0)}, h^{\xi,(0)}$ and per-scenario profits $f^{\xi,(0)}$. The expected initial profit is $f^{(0)} = \sum_\xi p_\xi\, f^{\xi,(0)}$, and the initial consensus values are computed node by node. The penalty parameter is auto-calibrated so that the average initial quadratic disagreement penalty equals a fixed fraction $\alpha = 0.1$ of $|f^{(0)}|$. Letting $\mathcal{Z}^\xi$ and $\mathcal{H}^\xi$ collect the $z$ and $h$ NA indices participating in scenario $\xi$, we set
\begin{align}
  \frac{\rho}{2} \cdot \frac{1}{S}
  \sum_{\xi=1}^{S} \Bigg[
  &\sum_{(u,s)\,\in\,\mathcal{Z}^\xi}
      \!\big(z^{\xi,(0)}_{u,s} - \bar{z}^{(0)}_{u,s}\big)^{\!2} \notag\\
  +\;&\sum_{(u,s,t)\,\in\,\mathcal{H}^\xi}
      \!\big(h^{\xi,(0)}_{u,s,t} - \bar{h}^{(0)}_{u,s,t}\big)^{\!2}
  \Bigg] \;=\; \alpha\,|f^{(0)}|.
\end{align}
and solve for $\rho$. The result is floored at a small positive value and capped at $0.05\,|f^{(0)}|$. The adaptive update during the loop is additionally bounded above by $10\,|f^{(0)}|$. Calibrating $\rho$ to a fraction of the expected objective is appropriate here because all binaries share the same $\{0,1\}$ scale, so a single scalar suffices.

\paragraph{Subproblem coefficient update.}
Applying the linearisation of the binary proximal term (Section~\ref{sec:ph_milp}) to each of the two blocks, the linear coefficient on each binary in the augmented objective is updated between iterations by
\begin{align}
  c^{\xi,z}_{u,s}   &\leftarrow -w^{\xi,z}_{u,s}   - \tfrac{\rho}{2}\big(1 - 2\bar{z}_{u,s,n(\xi)}\big),\\
  c^{\xi,h}_{u,s,t} &\leftarrow -w^{\xi,h}_{u,s,t} - \tfrac{\rho}{2}\big(1 - 2\bar{h}_{u,s,t,n(\xi)}\big),
\end{align}
where $n(\xi)$ denotes the relevant tree node for scenario $\xi$ at the corresponding stage. The constraint matrix is unchanged between iterations, only these coefficients are.

\paragraph{Consensus fixing.}
A binary NA variable is fixed at iteration $k$ if every scenario passing through its tree node returns the same integer value (within tolerance $\varepsilon_{\mathrm{bin}}=0.01$) for $\tau = 3$ consecutive iterations. Specifically, a per-variable unanimous count is incremented whenever all node scenarios agree, and reset to zero otherwise; the variable is fixed when the count reaches $\tau$. Its bounds are locked to that integer in every scenario MILP, and its multipliers are zeroed so the variable disappears from the augmented objective.

\paragraph{Slam trigger.}
Once $k \geq k_{\min} = 10$ iterations have elapsed, the slam is triggered when all three of the following hold simultaneously over a rolling window of $w_{\mathrm{slam}} = 6$ iterations: (i) the disagreement count has not fallen by more than $1\%$ relative to the first value in the window (\emph{disagree stagnated}); (ii) the maximum binary deviation $\max\text{Dev}$ has not fallen by more than $1\%$ relative to the first value in the window (\emph{dev stagnated}); and (iii) the current disagreement count is below $\delta_{\mathrm{slam}} = 0.5$ times the peak disagreement seen so far (\emph{made sufficient progress}). Conditions (i) and (ii) together detect a genuine stall; condition (iii) prevents the slam from firing prematurely before the algorithm has made meaningful progress.

\paragraph{Slam action (cycle-breaking fix).}
Once the slam is triggered, every unfixed binary NA variable whose consensus history $\bar{x}^{(j)}$ over the last $w_{\mathrm{slam}}$ iterations spans more than $\delta_{\mathrm{vol}} = 0.15$ between minimum and maximum is fixed to the integer nearest the mean of those values. Variables whose consensus has been stable are left for consensus fixing to handle. The oscillation-based trigger and action are our own design, inspired by the general slamming idea of \citet{watson2011progressive}.

\paragraph{Tail-stall escape.}
After the slam has been triggered, a separate tail-stall check runs every iteration. If the disagreement count is both at most $d_{\max} = 30$ and has been perfectly flat (unchanged) for $w_{\mathrm{tail}} = 5$ consecutive iterations, every remaining unfixed binary is force-fixed to the integer nearest its current $\bar{x}$, the convergence flag is set, and the loop exits. This prevents the algorithm from running indefinitely when a small handful of binaries refuse to settle after slamming.

\paragraph{Convergence.}
Phase~1 convergence is declared when there are no more disagreeing binary variables, following the idea of \emph{integer convergence} introduced by \citet{lokketangen1996progressive}.
This is checked after every iteration.

\paragraph{Adaptive penalty.}
Within the main loop $\rho$ is adapted from a rolling disagreement window
of length $w = 3$. When disagreement stagnates, a higher $\rho$ makes the
penalty term heavier relative to the original objective, incentivising each
scenario subproblem to move toward consensus even at the cost of a worse
individual objective value. Conversely, when disagreement is already
falling, relaxing $\rho$ prevents the penalty from overwhelming the
objective and allows the subproblems to optimise solution quality more
freely. Concretely, if the window's minimum disagreement has not fallen
below $80\%$ of the first value in the window, $\rho$ is multiplied by
$1.3$. If instead the most recent disagreement is below $80\%$ of the
first value, $\rho$ is divided by $1.3$ but never reduced below the
initial calibrated value. The asymmetry is deliberate: sustained stagnation over the window is required to trigger a penalty hike, while a single improving iteration is sufficient to reduce it.

\paragraph{Multiplier update.}
After each solve, the Lagrange multipliers are updated by the standard Rockafellar--Wets rule applied to both binary blocks (the node index $n(\xi)$ is implicit in the superscript notation):
\begin{align}
  w^{\xi,z}_{u,s}   &\leftarrow w^{\xi,z}_{u,s}   + \rho\,\big(z^\xi_{u,s} - \bar{z}_{u,s,n(\xi)}\big),\\
  w^{\xi,h}_{u,s,t} &\leftarrow w^{\xi,h}_{u,s,t} + \rho\,\big(h^\xi_{u,s,t} - \bar{h}_{u,s,t,n(\xi)}\big),
\end{align}
applied only to scenario--variable pairs where the scenario participates in the corresponding tree node. Multipliers for already-fixed variables are immediately zeroed out.

\paragraph{Consistency repair.}
A post-hoc repair pass corrects logical contradictions that can arise
when $z$ and $h$ are slammed independently, without the slam mechanism
checking their joint logical relationship. This is described as the
third contribution of B-PHA in Section~\ref{sec:lit_vs_own} and
formalised in Algorithm~\ref{alg:volfix}.

\paragraph{Pseudocode.}
The Phase~1 loop is split across Algorithms~\ref{alg:phase1}
and~\ref{alg:phase1b} for typesetting reasons and together they form a
single loop. The consensus-fixing subroutine is in
Algorithm~\ref{alg:consensus}. The slam trigger and action are in
Algorithm~\ref{alg:slam}. The post-processing consistency repair that
runs once before Phase~2 is in Algorithm~\ref{alg:volfix}. The
Phase~2 extensive-form LP resolve for stocking quantities is in
Algorithm~\ref{alg:phase2}.

\begin{algorithm}[htbp]
\caption{Phase 1 --- B-PHA main loop (Part~1: setup and inner loop steps)}
\label{alg:phase1}
\KwIn{Scenario MILPs $\{\mathcal{P}^\xi\}_{\xi=1}^{S}$;
      binary NA decisions: stocking timing $z$ and harvest timing $h$, per tree node; $K;
      $\alpha, \varepsilon_{\mathrm{bin}}, \tau, w, w_{\mathrm{slam}},
       k_{\min}, \delta_{\mathrm{slam}}, \delta_{\mathrm{vol}},
       w_{\mathrm{tail}}, d_{\max}$}
\KwOut{Node-wise production calendar $(\bar{z}, \bar{h})$
       --- continued in Algorithm~\ref{alg:phase1b}}
\BlankLine
\tcp{Step 0: initial solve and calibration}
\lForEach{$\xi$}{solve $\mathcal{P}^\xi$ (no penalty)
  $\to z^{\xi,(0)}, h^{\xi,(0)}, f^{\xi,(0)}$}
$f^{(0)}\leftarrow\sum_\xi p_\xi f^{\xi,(0)}$;\quad
compute node-wise averages $\bar{z}^{(0)}, \bar{h}^{(0)}$\;
calibrate $\rho$: initial penalty $=\alpha|f^{(0)}|$;\quad
$\rho\leftarrow\min(\rho,\,0.05|f^{(0)}|)$\;
$\rho_{\mathrm{cal}}\leftarrow\rho$;\quad
$\rho_{\max}\leftarrow\rho_{\mathrm{maxfrac}}\cdot|f^{(0)}|$\;
multipliers $\leftarrow 0$;\quad
\textsc{unanimous}$[i]\leftarrow 0$; \quad $k\leftarrow 0$\;
\textsc{fixed}$[i]\leftarrow\texttt{false}$ $\forall i$;\quad
\textsc{slam}$\leftarrow\texttt{false}$\;
\BlankLine
\While{$\neg{converged}$}{
   $k\leftarrow k + 1$\;
  \tcp{update augmented objective in every $\mathcal{P}^\xi$}
  \lForEach{unfixed binary $i$}{
    $c^{\xi,z}_{u,s}
      \leftarrow -w^{\xi,z}_{u,s}
        - \tfrac{\rho}{2}\!\left(1 - 2\,\bar{z}^{(k-1)}_{u,s,\,n(\xi)}\right)$\;
    $c^{\xi,h}_{u,s,t}
      \leftarrow -w^{\xi,h}_{u,s,t}
        - \tfrac{\rho}{2}\!\left(1 - 2\,\bar{h}^{(k-1)}_{u,s,t,\,n(\xi)}\right)$}
  \tcp{solve each scenario subproblem; update node consensus}
  \lForEach{$\xi$}{solve $\mathcal{P}^\xi$
    $\to z^{\xi,(k)}, h^{\xi,(k)}$}
  $\bar{z}^{(k)}_{u,s,n}\leftarrow
    \sum_{\xi\in n} p_{\xi|n}\,z^{\xi,(k)}_{u,s}$,\quad
  $\bar{h}^{(k)}_{u,s,t,n}\leftarrow
    \sum_{\xi\in n} p_{\xi|n}\,h^{\xi,(k)}_{u,s,t}$\;
  \BlankLine
  \tcp{consensus fixing (Algorithm~\ref{alg:consensus})}
  \textsc{ConsensusFix}($\bar{z}^{(k)}, \bar{h}^{(k)},
    \{z^{\xi,(k)}, h^{\xi,(k)}\}$)\;
  \BlankLine
  \tcp{slam trigger and action (Algorithm~\ref{alg:slam})}
  \textsc{SlamCheck}(\textsc{slam}, $k$, $k_{\min}$,
    $\mathrm{disagree\_hist}$, $\mathrm{maxdev\_hist}$,
    $w_{\mathrm{slam}}$, $\delta_{\mathrm{slam}}$, $\delta_{\mathrm{vol}}$,
    $\{\bar{z}^{(j)}, \bar{h}^{(j)}\}_{j\leq k}$)\;
  \BlankLine
  \tcp{--- continued in Algorithm~\ref{alg:phase1b} ---}}
\end{algorithm}

\begin{algorithm}[htbp]
\caption{Phase 1 --- B-PHA main loop (Part~2: continuation of Algorithm~\ref{alg:phase1})}
\label{alg:phase1b}
\BlankLine
\tcp{(inside the same for-loop, continuing from Algorithm~\ref{alg:phase1})}
\BlankLine
  \tcp{convergence check}
  $\mathrm{maxDev}\leftarrow\max\bigl(
    \max_{\xi,u,s}|z^{\xi,(k)}_{u,s}-\bar{z}^{(k)}_{u,s,n}|,\;
    \max_{\xi,u,s,t}|h^{\xi,(k)}_{u,s,t}-\bar{h}^{(k)}_{u,s,t,n}|
  \bigr)$ (unfixed, participating)\;
  $\mathrm{disagree}\leftarrow$ number of unfixed NA nodes $i$ where
    any scenario $\xi$ deviates from the node average by more than $\varepsilon_{\mathrm{bin}}$\;
  append $\mathrm{maxDev}$ to $\mathrm{maxdev\_hist}$, $\mathrm{disagree}$ to $\mathrm{disagree\_hist}$\;
  converged$\leftarrow(\mathrm{maxDev}\leq\varepsilon_{\mathrm{bin}})$\;
  \BlankLine
  \tcp{tail-stall escape}
  \If{\textnormal{not} converged \textbf{and} \textsc{slam}
      \textbf{and} $0<\mathrm{disagree}\leq d_{\max}$
      \textbf{and} $|\mathrm{disagree\_hist}|\geq w_{\mathrm{tail}}$}{
    \If{last $w_{\mathrm{tail}}$ entries of $\mathrm{disagree\_hist}$ all equal}{
      fix all remaining unfixed $z_{u,s,n}$ to $\mathrm{round}(\bar{z}^{(k)}_{u,s,n})$\;
      fix all remaining unfixed $h_{u,s,t,n}$ to $\mathrm{round}(\bar{h}^{(k)}_{u,s,t,n})$\;
      zero all remaining multipliers;\quad
      converged$\leftarrow\texttt{true}$\;
    }
  }
  \lIf{converged}{\textbf{break}}
  \BlankLine
  \tcp{multiplier update}
  $w^{\xi}_{z,u,s,n}\leftarrow w^{\xi}_{z,u,s,n}
    +\rho\,(z^{\xi,(k)}_{u,s}-\bar{z}^{(k)}_{u,s,n})$\quad
  for all participating \& unfixed stocking variables\:
  $w^{\xi}_{h,u,s,t,n}\leftarrow w^{\xi}_{h,u,s,t,n}
    +\rho\,(h^{\xi,(k)}_{u,s,t}-\bar{h}^{(k)}_{u,s,t,n})$\quad
  for all participating \& unfixed harvest variables\:
  \BlankLine
  \tcp{adaptive $\rho$, rolling window of length $w$}
  let $\mathrm{disagree\_window} \leftarrow \mathrm{disagree\_hist}[-w:]$\;
  \eIf{$\min(\mathrm{disagree\_window})>0.8\cdot \mathrm{disagree\_window}[0]$}{
    $\rho\leftarrow\min(\rho\cdot 1.3,\,\rho_{\max})$\;
  }{
    \lIf{$\mathrm{disagree\_window}[-1]<0.8\cdot \mathrm{disagree\_window}[0]$}{
      $\rho\leftarrow\max(\rho/1.3,\,\rho_{\mathrm{cal}})$}
  }
\BlankLine
\Return stocking calendar $\bar{z}$ and harvest calendar $\bar{h}$
from fixed and rounded node values\;
\end{algorithm}

\begin{algorithm}[htbp]
\caption{\textsc{ConsensusFix} --- agreement-based fixing on naturally settled nodes}
\label{alg:consensus}
\KwIn{Current per-scenario binary solutions
      $\{z^{\xi,(k)}, h^{\xi,(k)}\}$;
      current node averages $\bar{z}^{(k)}, \bar{h}^{(k)}$;
      tolerance $\varepsilon_{\mathrm{bin}}$; streak threshold $\tau$}
\BlankLine
\ForEach{unfixed stocking timing node $z_{u,s}$ at tree node $n$}{
  $v \leftarrow \mathrm{round}(z^{\xi_0,(k)}_{u,s})$ for any $\xi_0 \in n$\;
  \uIf{$|z^{\xi,(k)}_{u,s} - v| \leq \varepsilon_{\mathrm{bin}}$ for \textbf{all} $\xi \in n$}{
    \textsc{unanimous}$[z_{u,s,n}] \mathrel{+}= 1$\;
    \lIf{\textsc{unanimous}$[z_{u,s,n}] \geq \tau$}{
      fix $\bar{z}_{u,s,n} \leftarrow v$ in every $\mathcal{P}^\xi$ with $\xi \in n$;
      zero multipliers for $z_{u,s,n}$}
  }
  \lElse{\textsc{unanimous}$[z_{u,s,n}] \leftarrow 0$}
}
\BlankLine
\ForEach{unfixed harvest timing node $h_{u,s,t}$ at tree node $n$}{
  $v \leftarrow \mathrm{round}(h^{\xi_0,(k)}_{u,s,t})$ for any $\xi_0 \in n$\;
  \uIf{$|h^{\xi,(k)}_{u,s,t} - v| \leq \varepsilon_{\mathrm{bin}}$ for \textbf{all} $\xi \in n$}{
    \textsc{unanimous}$[h_{u,s,t,n}] \mathrel{+}= 1$\;
    \lIf{\textsc{unanimous}$[h_{u,s,t,n}] \geq \tau$}{
      fix $\bar{h}_{u,s,t,n} \leftarrow v$ in every $\mathcal{P}^\xi$ with $\xi \in n$;
      zero multipliers for $h_{u,s,t,n}$}
  }
  \lElse{\textsc{unanimous}$[h_{u,s,t,n}] \leftarrow 0$}
}
\end{algorithm}

\begin{algorithm}[htbp]
\caption{\textsc{SlamCheck} --- stall detection, oscillation fixing, and tail-stall escape}
\label{alg:slam}
\KwIn{Flag \textsc{slam}; iteration $k$; minimum iterations $k_{\min}$;
      history arrays $\mathrm{disagree\_hist}$, $\mathrm{maxdev\_hist}$;
      window $w_{\mathrm{slam}}$; stagnation threshold $\delta_{\mathrm{slam}}$;
      oscillation threshold $\delta_{\mathrm{vol}}$;
      consensus histories $\{\bar{z}^{(j)}, \bar{h}^{(j)}\}_{j \leq k}$}
\KwOut{Updated flag \textsc{slam}; newly fixed stocking/harvest timing nodes}
\BlankLine
\tcp{Slam trigger: detect stall once enough iterations have elapsed}
\If{$\neg$\textsc{slam} \textbf{and} $k \geq k_{\min}$}{
  let $\mathrm{disagree\_window} \leftarrow \mathrm{disagree\_hist}[-w_{\mathrm{slam}}:]$
  (last $w_{\mathrm{slam}}$-iterations' disagreement counts)\;
  let $\mathrm{maxdev\_window} \leftarrow \mathrm{maxdev\_hist}[-w_{\mathrm{slam}}:]$
  (last $w_{\mathrm{slam}}$-iterations' max-deviation values)\;
  stagnant $\leftarrow \min(\mathrm{disagree\_window})
    \geq (1 - \delta_{\mathrm{slam}}) \cdot \mathrm{disagree\_window}[0]$
    \textbf{and} $\mathrm{maxdev\_window}$ is flat\;
  progress $\leftarrow \mathrm{disagree\_window}[-1]
    < 0.5 \cdot \max(\mathrm{disagree\_hist})$\;
  \lIf{stagnant \textbf{and} progress}{\textsc{slam} $\leftarrow \texttt{true}$}
}
\BlankLine
\tcp{Slam action: fix each binary type whose consensus has oscillated}
\If{\textsc{slam}}{
  \ForEach{unfixed stocking timing node $z_{u,s}$ at tree node $n$}{
    \lIf{$\max_j \bar{z}^{(j)}_{u,s,n} - \min_j \bar{z}^{(j)}_{u,s,n} > \delta_{\mathrm{vol}}$}{
      fix $z_{u,s,n} \leftarrow \mathrm{round}(\mathrm{mean}_j \bar{z}^{(j)}_{u,s,n})$;
      zero multipliers for this node}
  }
  \ForEach{unfixed harvest timing node $h_{u,s,t}$ at tree node $n$}{
    \lIf{$\max_j \bar{h}^{(j)}_{u,s,t,n} - \min_j \bar{h}^{(j)}_{u,s,t,n} > \delta_{\mathrm{vol}}$}{
      fix $h_{u,s,t,n} \leftarrow \mathrm{round}(\mathrm{mean}_j \bar{h}^{(j)}_{u,s,t,n})$;
      zero multipliers for this node}
  }
}
\end{algorithm}

\begin{algorithm}[htbp]
\caption{\textsc{ConsistencyRepair} --- post-hoc $z$--$h$ repair before Phase~2}
\label{alg:volfix}
\KwIn{Binary production calendar $(\bar{z}, \bar{h})$
      from the B-PHA loop; fixed flags \textsc{fixed}$[i]$,
      force-fix flags \textsc{slammed}$[i]$ (true if fixed by slam or
      tail-stall escape rather than natural consensus fixing),
      fix values \textsc{fix\_val}$[i]$, and final consensus
      $\bar{x}^{(k)}$ for all binary NA variables;
      NA name-to-index map; tree node assignments}
\KwOut{Repaired production calendar where every force-fixed $h_{u,s,t}=1$
       has a corresponding $z_{u,s}=1$. %and every force-fixed $z_{u,s}=1$ has at least one corresponding $h_{u,s,t}=1$ at the governing tree node
       }
\BlankLine
\tcp{Assemble integer consensus from PH output}
\ForEach{binary NA variable $i$}{
  \lIf{\textsc{fixed}$[i]$}{$\bar{x}_i \leftarrow \textsc{fix\_val}[i]$}
  \lElse{$\bar{x}_i \leftarrow \mathrm{round}(\bar{x}^{(k)}_i)$}
}
\BlankLine
\tcp{Forward check: cancel force-fixed harvests whose cohort was never stocked}
\ForEach{harvest timing node $h_{u,s,t}$ at tree node $n_h$
    with $\bar{h}_{u,s,t,n_h} = 1$ \textbf{and} \textsc{slammed}$[h_{u,s,t,n_h}]$}{
  derive $n_z \leftarrow$ tree node governing stocking month $s$
  (ancestor of $n_h$ if $s$ belongs to an earlier stage)\;
  \lIf{$\bar{z}_{u,s,n_z} < 0.5$}{
    $\bar{h}_{u,s,t,n_h} \leftarrow 0$
    \tcp*[r]{cohort never stocked --- cancel harvest}}
}
\BlankLine
\tcp{Backward check: cancel force-fixed stocking whose cohort has no harvest}
\ForEach{stocking timing node $z_{u,s}$ at tree node $n_z$
    with $\bar{z}_{u,s,n_z} = 1$ \textbf{and} \textsc{slammed}$[z_{u,s,n_z}]$}{
  \lIf{$\bar{h}_{u,s,t,n} = 0$ for \textbf{all} harvest nodes $h_{u,s,t}$
      at tree nodes $n \supseteq n_z$}{
    $\bar{z}_{u,s,n_z} \leftarrow 0$
    \tcp*[r]{no harvest available --- cancel stocking}}
}
\BlankLine
\tcp{Apply repaired production calendar as fixed bounds to every scenario MILP}
\ForEach{$\xi = 1,\ldots,S$}{
  \ForEach{stocking timing node $z_{u,s}$ at tree node $n$ participating in $\mathcal{P}^\xi$}{
    set $\mathrm{LB} = \mathrm{UB} = \bar{z}_{u,s,n}$ in $\mathcal{P}^\xi$\;
  }
  \ForEach{harvest timing node $h_{u,s,t}$ at tree node $n$ participating in $\mathcal{P}^\xi$}{
    set $\mathrm{LB} = \mathrm{UB} = \bar{h}_{u,s,t,n}$ in $\mathcal{P}^\xi$\;
  }
  restore penalty-free objective in $\mathcal{P}^\xi$\;
}
\end{algorithm}
%-----------------------------------------------------------
\subsubsection{Phase 2 --- Stochastic LP resolve for stocking quantities}
\label{sec:lpresolve}
%-----------------------------------------------------------

Phase~1 returns the production calendar $(\bar{z}, \bar{h})$, with timing identical across scenarios that share a tree node. Once these binaries are fixed, the per-scenario MILP $\mathcal{P}^\xi$ from Section~\ref{sec:compact} collapses to a linear program in the continuous stocking quantity $q$ together with the auxiliary continuous variables biomass $B^\xi$ and harvested biomass $H^\xi_{\mathrm{biom}}$.

% Kommentere ut hele avsnittet under?
% \paragraph{Why this can be solved in extensive form.}
% Without integer variables, the problem is convex, and its deterministic equivalent, a single LP combining all 81 scenarios with non-anticipativity constraints on $q$, is tractable. We solve it directly. Non-anticipativity on $q$ is enforced \emph{structurally} rather than through equality constraints: each $q$ variable is indexed by tree node $n$, written $q_{u,s,n}$, so that scenarios sharing $n$ share the same stocking quantity by construction. This eliminates the coupling that made the original SMILP hard, and the resulting problem is a standard convex LP solvable in seconds with a single barrier-method call. 

\paragraph{Per-scenario evaluation.}
Once the extensive-form LP returns optimal stocking quantities
$q^*_{u,s,n}$, each scenario is evaluated individually with the full
plan $(\bar{z}, \bar{h}, q^*)$ fixed to recover the per-scenario
objective $f^\xi$. These per-scenario values are needed for VSS and
EVPI reporting in Sections~\ref{sec:Determinisitic_compared_Stochastic}
and~\ref{sec:sp_vs_ef}, and for the scenario-level diagnostics
discussed there.

\paragraph{Pseudocode.}
The full Phase~2 procedure is shown in Algorithm~\ref{alg:phase2}.

\begin{algorithm}[htbp]
\caption{Phase 2 --- Extensive-form stochastic LP for $q$}
\label{alg:phase2}
\KwIn{Repaired node-wise production calendar $(\bar{z},\bar{h})$
      from Algorithm~\ref{alg:volfix};
      scenario MILPs $\{\mathcal{P}^\xi\}_{\xi=1}^{S}$}
\KwOut{Node-wise stocking quantities $q^*$ and per-scenario profits
      $\{f^\xi\}_{\xi \in \Xi}$}
\BlankLine
\tcp{Build the extensive-form LP}
initialise an empty LP $\mathcal{L}$\;
\ForEach{$\xi = 1,\ldots,S$}{
  fix stocking timing $z^\xi_{u,s} \leftarrow \bar{z}_{u,s,n(\xi)}$
  and harvest timing $h^\xi_{u,s,t} \leftarrow \bar{h}_{u,s,t,n(\xi)}$
  in $\mathcal{P}^\xi$\;
}
add node-indexed stocking quantity variables $q_{u,s,n}$
  for every cage $u$, stocking month $s$, and tree node $n$\;
\ForEach{stage $k \in \{0,1,2,3\}$}{
  \ForEach{pair of scenarios $\xi, \xi'$ sharing node $n$ at stage $k$}{
    enforce $q^\xi_{u,s} = q^{\xi'}_{u,s}$ for all $u, s \in T_k$
    \tcp*[r]{hard NACs on $q$ by stage}
  }
}
% add elastic slack variables for MAB and density constraints,
%   priced at $c^{\mathrm{pen}}$, to guarantee feasibility\;
\ForEach{$\xi = 1,\ldots,S$}{
  add scenario-$\xi$ biomass-recursion, density, location-MAB,
  and regional-MAB constraints, expressed in terms of the node-shared $q_{u,s,n}$\;
  add scenario-$\xi$ revenue, harvesting cost, smolt cost, and terminal-value terms,
  weighted by $p_\xi$, to the objective\;
}
\BlankLine
solve $\mathcal{L}$ as a single extensive-form LP $\to q^*$\;
\BlankLine
\tcp{Per-scenario evaluation: fix full plan and evaluate profit for each scenario}
\ForEach{$\xi = 1,\ldots,S$}{
  fix $z^\xi_{u,s}, h^\xi_{u,s,t}, q^\xi_{u,s}$ in $\mathcal{P}^\xi$
  to $(\bar{z}_{u,s,n(\xi)}, \bar{h}_{u,s,t,n(\xi)}, q^*_{u,s,n(\xi)})$\;
  evaluate per-scenario profit $f^\xi$ by solving $\mathcal{P}^\xi$\;
}
\Return $q^*, \{f^\xi\}_{\xi \in \Xi}$\;
\end{algorithm}

\newpage

    
\section{Experimental Evaluation}
\label{experiment}

This section evaluates the proposed Progressive Hedging heuristic against
the deterministic equivalent (DE), the expected value of the expected value plan (EEV), and the
wait-and-see bound (WS) across six fleet sizes ranging from 4 to 60 cages.
EVPI and VSS are computed relative to the 81-scenario tree described
below, not relative to the historical distribution of Norwegian sea
temperatures. The point of the tree is not to reproduce the empirical
distribution but to incorporate uncertainty by spanning a credible range
of growth and mortality outcomes; the absolute EVPI and VSS numbers are
therefore conditional on the scenario generation choices, while the
qualitative finding that flexible stochastic planning has structural
value whenever there is temperature uncertainty 
% creates materially different
% trajectories does not depend on the exact probability weights.
% Out-of-sample experiments on historical data would speak more directly
% to the realised gain of stochastic modelling, and we leave this to
% producers and future researchers.

% All experiments were run on a base-model Apple MacBook Air (M2 chip, 8 GB
% unified memory, 256 GB SSD) using Gurobi as the underlying MILP solver,
% which is modest hardware for this industry. The wall-clock and memory
% results below should be interpreted with this baseline in mind.

\subsection{Dataset and problem-size scaling}
\label{sec:data}

\paragraph{The instance.}
The dataset is calibrated against operational data and parameters supplied
by Scale Aquaculture, a professional salmon farming service firm. Biological
growth and survival parameters are drawn from the literature and assumed
adequate for the purposes of this thesis. The instance defines a planning
horizon of 60 months (five years) and a fleet of 60 cages partitioned
across six coastal locations (10, 12, 8, 10, 8, and 12 cages
respectively). Each cage is characterised by its physical volume
(30{,}000 m\textsuperscript{3} or 35{,}000 m\textsuperscript{3}) and its
location. A subset of 18 cages start the horizon already stocked at given
headcounts (typically between 145{,}000 and 190{,}000 fish) and average
weights ranging from approximately 500 g for recently stocked cohorts to
5{,}700 g for cohorts near harvest weight. The remaining 42 cages begin
empty and must be stocked during the horizon if they are to contribute
to the plan. The six locations have mandatory coordinated fallow periods
every 24 months, with the first fallow month at month 12, 13, 14, 15,
16, and 18 respectively (the staggering across locations is taken from
the operational data and is not a regular sequence). Per-location MAB
limits are 4{,}000, 5{,}600, 2{,}500, 4{,}300, 2{,}600, and 5{,}500
tonnes respectively, with a regional MAB cap of 35{,}000 tonnes
governing the aggregate biomass across all locations at any month. These
limits are enforced softly through slack variables with a 1{,}000 NOK
per kg penalty, strongly discouraging violations in the
optimal solution. Per-cage stocking density is bounded at
25 kg/m\textsuperscript{3}.

\paragraph{Experimental fleet sizes.}
The experiments use six fleet sizes (4, 8, 10, 12, 30, and 60 cages) selected as subsets of the full 60-cage fleet. The smaller sizes restrict to a single location, while the larger sizes activate cages across all six locations to exercise the regional MAB constraint as well as the per-location limits.

\paragraph{Scenario tree.}
Biological uncertainty is represented through a four-stage scenario tree
covering temperature realisations across the horizon. Each stage spans 15
months (months 0 to 14, 15 to 29, 30 to 44, and 45 to 59), and each stage
branches into three equiprobable temperature outcomes (cold, normal,
warm) with monthly profiles offset by $\pm 2$\textdegree C from the
normal seasonal pattern. This produces $3^4 = 81$ scenarios, each with
probability $1/81$. Survival rates are coupled to the temperature
realisation: a monthly survival rate of $(1 - 0.0002)^{30}$ applies under
normal conditions, dropping to $(1 - 0.001)^{30}$ when a warm stage
follows another warm stage, capturing the increased sea lice pressure
that warm-water sequences can produce in practice. Temperature also
drives fish growth through standard biological growth functions
evaluated against the scenario-specific temperature profile, producing
scenario-specific survival-adjusted biomass-per-fish curves that enter
both the constraints and the objective.

\paragraph{Why DE scales poorly while SP does not.}
DE assembles all 81 scenarios into a single monolithic model and enforces
non-anticipativity through hard equality constraints, so every cage in every
scenario contributes its own stocking and harvest binaries that must be
coordinated simultaneously. The branch-and-bound search space therefore grows
with the total number of binary decisions across all scenarios at once, and
adding cages multiplies this burden because each new cage couples into all 81
scenarios together.

The SP heuristic decomposes along the scenario axis, solving 81 single-scenario
subproblems that are coordinated through penalty terms in the objective rather
than hard linking constraints. Each subproblem has the same combinatorial
structure as a single deterministic IP, and crucially the penalty is applied
only to the binary timing decisions, not to the continuous stocking
quantities, so no additional coupling is introduced. Once binary consensus is
reached across scenarios, the remaining recourse problem reduces to a pure
linear program solved in extensive form. This separation of the scenario axis from the combinatorial axis is the
structural reason the SP remains tractable at scales where the DE does not.

\subsection{Deterministic versus stochastic planning}
\label{sec:Determinisitic_compared_Stochastic}

We first compare the EEV, DE, and WS solvers at scales where DE remains
tractable; the SP heuristic is introduced in §\ref{sec:sp_vs_ef}.
Table~\ref{tab:model_stats_cages} reports results for 4, 8, and 12 cages.
The 4-cage and 8-cage experiments use a 2\% MIP gap tolerance across all
solvers, while the 12-cage experiment uses a 3.5\% tolerance to keep DE
within a manageable runtime budget.

\begin{table}[htbp]
\centering
\caption{Model statistics and performance by cage configuration}
\label{tab:model_stats_cages}
\footnotesize
\begin{tabularx}{\linewidth}{lRRRRRRRRR}
\toprule
& \multicolumn{3}{c}{\textbf{4 Cages}} & \multicolumn{3}{c}{\textbf{8 Cages}} & \multicolumn{3}{c}{\textbf{12 Cages}} \\
\cmidrule(lr){2-4} \cmidrule(lr){5-7} \cmidrule(lr){8-10}
\textbf{Metric} & \textbf{EEV} & \textbf{DE} & \textbf{WS}
               & \textbf{EEV} & \textbf{DE} & \textbf{WS}
               & \textbf{EEV} & \textbf{DE} & \textbf{WS} \\
\midrule
Scenarios              & 54/81\textsuperscript{*} & 81 & 81
                       & 54/81\textsuperscript{*} & 81 & 81
                       & 54/81\textsuperscript{*} & 81 & 81 \\
E[obj] (MNOK)          & 325.9 & 360.4 & 426.6
                       & 493.1 & 549.7 & 644.2
                       & 692.2 & 737.9 & 847.0 \\
MIP gap tolerance (\%) & 2     & 2     & 2
                       & 2     & 2     & 2
                       & 3.5   & 3.5   & 3.5 \\
Wall clock time (s)    & 31.6  & 55.5  & 50.0
                       & 58.7  & 248.0 & 113.1
                       & 179.2 & 1{,}652 & 372.4 \\
\midrule
VSS (MNOK)  & \multicolumn{3}{r}{$\text{DE} - \text{EEV} = 34.5$}
            & \multicolumn{3}{r}{$\text{DE} - \text{EEV} = 56.6$}
            & \multicolumn{3}{r}{$\text{DE} - \text{EEV} = 45.7$} \\
EVPI (MNOK) & \multicolumn{3}{r}{$\text{WS} - \text{DE} = 66.2$}
            & \multicolumn{3}{r}{$\text{WS} - \text{DE} = 94.5$}
            & \multicolumn{3}{r}{$\text{WS} - \text{DE} = 109.1$} \\
\bottomrule
\end{tabularx}
\par\vspace{2pt}
\footnotesize\textsuperscript{*}EEV expectation taken over the 54 feasible scenarios (probability mass 0.667). 27 scenarios were infeasible under the EV plan at every scale.
\end{table}

Two patterns appear immediately. First, DE wall-clock time accelerates
sharply with fleet size, from 56 seconds at 4 cages to 248 seconds at 8
cages and 1{,}652 seconds at 12 cages, roughly fourfold and thirtyfold
increases respectively, while the EEV and WS times grow far more
modestly. This is the empirical signature of the combinatorial growth
analysed in §\ref{sec:data} and foreshadows the intractability that
emerges at 30 and 60 cages in §\ref{sec:sp_vs_ef}. Second, the EV plan is infeasible 
in exactly 27 of 81 scenarios at every
scale, a ratio that persists across all fleet sizes we tested (see also
Tables~\ref{tab:model_stats_10cages}, \ref{tab:model_stats_30cages}, and
\ref{tab:model_stats_large}). 
The infeasibilities arise when 
%from how uncertainty is discretised in this model: the coarse three-branch temperature tree produces scenario realisations that can deviate enough from the expected path that 
the expected-value (EV) optimal plan becomes biologically infeasible in the cold-temperature branches. EEV is therefore reported as a conditional expectation over the 54 feasible scenarios. That means the realised profit under every feasible scenario is weighted by the factor $1/54$ instead of $1/81$ when looking at the EEV. This is a slightly biased comparison and the reported VSS values should be read as such, but, we found it to be a more meaningful comparison than one where a third of the scenarios returns a 0 NOK profit under the EV-plan due to slight infeasibilities (discussed in \ref{sec:sensitivity_params}). We further speculate that the reported EEV could be a seen as a upper bound as it does not contain realised profit under scenarios starting with a cold stage, which of course is a deviation from the normal temperature scenario. And, since the EV-plan only plans for normal temperatures, then these 27 scenarios which were deemed infeasible would probably contribute a less than average profit if they were to be feasible. We therefore claim that the EEV reported here is upwards biased. This is of course speculative as one cannot know the true value of an infeasible plan. But what can be said is that the deterministic plan does not merely produce worse outcomes in expectation, it produces no executable plan at all in roughly a third of futures, a robustness failure that the conditional VSS metric probably understates.


The conditional VSS grows from 34.5~MNOK at 4 cages to 56.6~MNOK at 8
cages, then dips to 45.7~MNOK at 12 cages. The non-monotonic behaviour
at 12 cages is potentially partly an artefact of the looser 3.5\% MIP
gap tolerance used at that scale, which can inflate the gap between the
DE incumbent and the true optimum and therefore might compress the
apparent VSS. Another contributing factor may be that all three fleet
sizes draw cages from the same single location, meaning additional cages
share the same fixed per-location $\mathrm{MAB}_\ell$ ceiling that existed for fewer cages. This could limit the marginal productive capacity each
new cage adds and thereby dampen the growth in VSS. EVPI in contrast
grows monotonically from 66.2 to 94.5 to 109.1~MNOK. 
% The contrast
% between the two trajectories is itself informative: the value of perfect
% information scales smoothly with fleet size while the measured value of stochastic programming is shaped by the conservative nature of making a feasible plan for an uncertain future. % solver tolerance, suggesting that VSS estimates at looser tolerances should be read as lower bounds on the true value.

\subsection{The SP heuristic against the extensive form}
\label{sec:sp_vs_ef}

We now turn to the central comparison of the thesis: how the SP heuristic
performs against the extensive-form DE across a range of practically
relevant fleet sizes.

\begin{table}[htbp]
\centering
\caption{Model statistics and performance, 10 cages}
\label{tab:model_stats_10cages}
\footnotesize
\begin{tabularx}{\linewidth}{lRRRR}
\toprule
\textbf{Metric} & \textbf{EEV} & \textbf{SP} & \textbf{DE} & \textbf{WS} \\
\midrule
Scenarios              & 54/81\textsuperscript{*} & 81 & 81 & 81 \\
E[obj] (MNOK)          & 634.8       & 661.4         & 691.8       & 789.5 \\
MIP gap tolerance (\%) & 2           & 2             & 2           & 2 \\
Wall clock time (s)    & 62.1        & 386.0         & 735.4       & 153.5 \\
\midrule
VSS$_{\text{SP}}$ (MNOK)  & \multicolumn{4}{r}{$\text{SP} - \text{EEV} = 26.6$} \\
VSS$_{\text{DE}}$ (MNOK)  & \multicolumn{4}{r}{$\text{DE} - \text{EEV} = 57.0$} \\
EVPI$_{\text{SP}}$ (MNOK) & \multicolumn{4}{r}{$\text{WS} - \text{SP} = 128.1$} \\
EVPI$_{\text{DE}}$ (MNOK) & \multicolumn{4}{r}{$\text{WS} - \text{DE} = 97.7$} \\
\bottomrule
\end{tabularx}
\par\vspace{2pt}
\footnotesize\textsuperscript{*}EEV expectation taken over the 54 feasible scenarios (probability mass 0.667). 27 scenarios were infeasible under the EV plan.
\end{table}

\begin{table}[htbp]
\centering
\caption{Model statistics and performance, 30 cages}
\label{tab:model_stats_30cages}
\footnotesize
\begin{tabularx}{\linewidth}{lRRRR}
\toprule
\textbf{Metric} & \textbf{EEV} & \textbf{SP} & \textbf{DE} & \textbf{WS} \\
\midrule
Scenarios              & 54/81\textsuperscript{*} & 81 & 81 & 81 \\
E[obj] (MNOK)          & 1{,}357.8     & 1{,}576.4   & 1{,}658.0\textsuperscript{\dag} & 1{,}920.6 \\
MIP gap tolerance (\%) & 3             & 3           & 3\textsuperscript{\dag}         & 3 \\
Wall clock time (s)    & 290.9         & 1{,}186.5   & 20{,}609                        & 747.5 \\
\midrule
VSS$_{\text{SP}}$ (MNOK)  & \multicolumn{4}{r}{$\text{SP} - \text{EEV} = 218.6$} \\
VSS$_{\text{DE}}$ (MNOK)  & \multicolumn{4}{r}{$\text{DE} - \text{EEV} = 300.2$} \\
EVPI$_{\text{SP}}$ (MNOK) & \multicolumn{4}{r}{$\text{WS} - \text{SP} = 344.2$} \\
EVPI$_{\text{DE}}$ (MNOK) & \multicolumn{4}{r}{$\text{WS} - \text{DE} = 262.6$} \\
\bottomrule
\end{tabularx}
\par\vspace{2pt}
\footnotesize\textsuperscript{*}EEV expectation taken over the 54 feasible scenarios (probability mass 0.667). 27 scenarios were infeasible under the EV plan.\\
\textsuperscript{\dag}DE terminated at 3.03\% gap, slightly above the 3\% tolerance.
\end{table}

\begin{table}[htbp]
\centering
\caption{Model statistics and performance, 60 cages (DE intractable)}
\label{tab:model_stats_large}
\footnotesize
\begin{tabularx}{\linewidth}{lRRRR}
\toprule
\textbf{Metric} & \textbf{EEV} & \textbf{SP} & \textbf{DE} & \textbf{WS}\\
\midrule
Scenarios              & 54/81\textsuperscript{*} & 81          & 81                      & 81 \\
E[obj] (MNOK)          & 2{,}383.5                & 3{,}235.5   & \multicolumn{1}{c}{OOM} & 3{,}870.6 \\
MIP gap tolerance (\%) & 2                        & 2           & \multicolumn{1}{c}{OOM} & 2 \\
Wall clock time (s)    & 509.5                    & 3{,}420.5   & \multicolumn{1}{c}{OOM} & 2{,}026.5 \\
\midrule
VSS$_{\text{SP}}$ (MNOK)  & \multicolumn{4}{r}{$\text{SP} - \text{EEV} = 852.0$} \\
EVPI$_{\text{SP}}$ (MNOK) & \multicolumn{4}{r}{$\text{WS} - \text{SP} = 635.1$} \\
\bottomrule
\end{tabularx}
\par\vspace{2pt}
\footnotesize\textsuperscript{*}EEV expectation taken over the 54 feasible scenarios (probability mass 0.667). 27 scenarios were infeasible under the EV plan.\\
\textit{Note:} OOM = out-of-memory. The DE extensive form exceeds the 8 GB of unified memory on the test machine at this scale (961{,}431 continuous and 866{,}111 integer variables, 53{,}226 SOS constraints).
\end{table}

Three findings emerge from this comparison.

\emph{The SP recovers solutions close to the extensive-form optimum.}
Where DE is tractable, the SP and DE objectives agree to within roughly
five percent. At 10 cages the gap is 4.4\% (661.4 versus 691.8 MNOK), and
at 30 cages it is 4.9\% (1{,}576.4 versus 1{,}658.0 MNOK). The binary consensus enforced by the SP-heuristic recovers solutions close to what a fully exact extensive-form, DE, solver would find. The VSS estimates for the SP and DE runs reported in table \ref{tab:model_stats_10cages} and \ref{tab:model_stats_30cages} show the value of stochastic programming from below
and above respectively: SP is a feasible non-anticipative solution and
DE is essentially exact, so arguably the true optimal VSS for the SP-algorithm lies between them.

\emph{The SP scales while DE does not.} Wall-clock time for the SP grows
from 386 to 1{,}187 to 3{,}421 seconds across 10, 30, and 60 cages,
scaling near-linearly in the number of cages (the second jump is mildly
super-linear, but only mildly). DE grows from 735 seconds at 10 cages to
20{,}609 seconds at 30 cages, a 28$\times$ increase for a 3$\times$
increase in fleet size, and at 60 cages the extensive form exceeds the
8 GB unified memory budget of the test machine entirely. The 17$\times$
SP-over-DE speedup at 30 cages and the complete intractability of DE at
60 cages establish the SP heuristic as the only method in this
comparison that produces a fully feasible non-anticipative plan with
near-optimal objective value at problem sizes that match real-world
operations on commodity hardware.

\emph{The value of resolving uncertainty grows with fleet size.} The
conditional VSS rises from 26.6 MNOK at 10 cages to 218.6 MNOK at 30
cages and 852.0 MNOK at 60 cages (using the SP), and EVPI grows in
parallel from 128.1 to 344.2 to 635.1 MNOK. Both reflect the intuition
that larger fleets expose more independent stocking and harvest
decisions to the same biological uncertainty, so the deterministic plan
accumulates more committed errors that the stochastic plan can hedge
against. 
% The growing EVPI shows that even the optimal non-anticipative plan
% leaves residual value relative to perfect information, reflecting the
% fundamental cost of having to commit to early-stage decisions before
% the temperature realisations are observed.
\newpage
\begin{figure}
    \centering
    \includegraphics[width=1\linewidth]{bilder/biomassplot.png}
    \caption{Total regional biomass trajectories under all 81 scenarios
    for the 60-cage instance solved by the SP heuristic. The thick blue
    line is the expectation across scenarios; thin grey lines are
    individual scenario trajectories. The plan stays comfortably below
    the 35{,}000-tonne regional MAB ceiling in every scenario, and
    scenario spread widens as growth uncertainty compounds across stages.}
    \label{fig:biomass_60cages_sp}
\end{figure}
\begin{figure}
    \centering
    \includegraphics[width=1.0\linewidth]{bilder/plan_full_dataset.png}
    \caption{Stocking and harvest plan along the all-normal path of the scenario tree for the 60-cage instance solved by the SP heuristic. Each colour represents one of the six locations; green triangles mark stocking events and orange triangles mark harvest events.}
    \label{fig:plan_normal_path_sp}
\end{figure}


\subsection{Scaling limits of the SP heuristic}
\label{sec:sp_70_cages}

To probe the upper limit of what the SP heuristic can handle on the test
hardware, we extended the instance to 70 cages by duplicating a subset of
the existing cage data into an additional location. The extension is synthetic, so the
absolute objective value should be read as a measure of computational
reach rather than a forecast for any real fleet. Table~\ref{tab:sp_70cages}
reports the result.

\begin{table}[htbp]
\centering
\caption{SP performance at 70 cages (synthetic extension)}
\label{tab:sp_70cages}
\footnotesize
\begin{tabularx}{\linewidth}{lR}
\toprule
\textbf{Metric} & \textbf{SP} \\
\midrule
Scenarios              & 81/81 feasible \\
E[obj] (MNOK)          & 3{,}881.1 \\
MIP gap tolerance (\%) & 2 \\
Wall clock time (s)    & 3{,}820.4 \\
\bottomrule
\end{tabularx}
\end{table}

The SP completes in 3{,}820 seconds with all 81 scenarios feasible at the
2\% gap tolerance, only 12\% slower than the 60-cage run (3{,}421
seconds) for a 17\% increase in fleet size. The scaling between 60 and 70
cages is therefore sub-linear, slightly better than the near-linear
pattern observed across 10, 30, and 60 cages. This strengthens the
tractability claim made in §\ref{sec:sp_vs_ef}. The SP not only solves
problems where DE went out-of-memory, it continues to scale past that point. At 80 cages and above, however, the SP itself exceeds
the 8 GB unified memory budget during execution, so the empirical
ceiling on this hardware sits between 70 and 80 cages.

The bottleneck appears to be the cumulative memory footprint of holding
81 scenario subproblems and their solver state simultaneously, rather
than any algorithmic intractability. Per-iteration solve times remain
comfortably within budget at 70 cages, and the failure mode at 80 cages
and beyond is memory exhaustion rather than runtime blow-up. We expect
that a workstation with more RAM, combined with parameter tuning of the
B-PHA penalty schedule and Gurobi-side memory settings, would push the ceiling further. 
% Sequential or batched subproblem solving, which trades wall-clock time for memory
% headroom, is also a natural next step for an implementation targeting
% larger fleets on commodity hardware.

\subsection{Sensitivity to temperature variance}
\label{sec:low_variance}

To examine how the value of stochastic programming responds to the
magnitude of uncertainty, we re-run the 10-cage experiment with a
narrower temperature spread between the cold, normal, and warm scenario
branches. The normal profile is unchanged, but the cold and warm profiles
are pulled in by 1\textdegree C from normal, halving the cold-warm spread
from 4\textdegree C to 2\textdegree C. The reduced profiles are shown in
Table~\ref{tab:reduced_temps}.

\begin{table}[htbp]
\centering
\caption{Monthly temperature profiles (\textdegree C) under the reduced-variance scenario tree. The same 12-month seasonal pattern is tiled across the 60-month horizon.}
\label{tab:reduced_temps}
\footnotesize
\begin{tabularx}{\linewidth}{lRRRRRRRRRRRR}
\toprule
\textbf{Branch} & \textbf{Jan} & \textbf{Feb} & \textbf{Mar} & \textbf{Apr} & \textbf{May} & \textbf{Jun} & \textbf{Jul} & \textbf{Aug} & \textbf{Sep} & \textbf{Oct} & \textbf{Nov} & \textbf{Dec} \\
\midrule
Cold     & 4.0 & 4.0 & 4.0 & 5.0 & 8.0  & 11.0 & 13.0 & 15.5 & 14.5 & 12.0 & 9.0  & 6.5 \\
Normal  & 5.0 & 5.0 & 5.0 & 6.0 & 9.0  & 12.0 & 14.0 & 16.5 & 15.5 & 13.0 & 10.0 & 7.5 \\
Warm    & 6.0 & 6.0 & 6.0 & 7.0 & 10.0 & 13.0 & 15.0 & 17.5 & 16.5 & 14.0 & 11.0 & 8.5 \\
\bottomrule
\end{tabularx}
\end{table}

The scenario tree structure (four stages, three equiprobable branches per
stage, 81 scenarios) is otherwise identical to the main experiments, as
are the survival rate coupling and all other instance parameters. All
four solvers run with a 2\% MIP gap tolerance.
Table~\ref{tab:model_stats_low_variance} reports the results.

\begin{table}[htbp]
\centering
\caption{Model statistics and performance, 10 cages with reduced temperature variance}
\label{tab:model_stats_low_variance}
\footnotesize
\begin{tabularx}{\linewidth}{lRRRR}
\toprule
\textbf{Metric} & \textbf{EEV} & \textbf{SP} & \textbf{DE} & \textbf{WS} \\
\midrule
Scenarios              & 81/81 feasible & 81 & 81 & 81 \\
E[obj] (MNOK)          & 703.0          & 735.5         & 753.0       & 809.3 \\
MIP gap tolerance (\%) & 2              & 2             & 2           & 2 \\
Wall clock time (s)    & 73.9           & 346.5         & 1{,}945.8   & 298.7 \\
\midrule
VSS$_{\text{SP}}$ (MNOK)  & \multicolumn{4}{r}{$\text{SP} - \text{EEV} = 32.5$} \\
VSS$_{\text{DE}}$ (MNOK)  & \multicolumn{4}{r}{$\text{DE} - \text{EEV} = 50.0$} \\
EVPI$_{\text{SP}}$ (MNOK) & \multicolumn{4}{r}{$\text{WS} - \text{SP} = 73.8$} \\
EVPI$_{\text{DE}}$ (MNOK) & \multicolumn{4}{r}{$\text{WS} - \text{DE} = 56.3$} \\
\bottomrule
\end{tabularx}
\end{table}

The clearest finding is that the EV plan is now feasible in all 81
scenarios, in contrast to the 27/81 infeasibility ratio observed under
the full-variance tree throughout this thesis. The underlying mechanism
is that under the original $\pm 2$\textdegree C spread, the EV plan optimises against an expected growth trajectory fast enough
to accommodate an additional production cycle within the horizon, with
stocking dates pushed late enough that under the worst cold-stage
realisations the fish cannot reach the minimum harvest weight required
by the well-boat and slaughter constraints before the mandatory fallow
periods. Compressing the cold branch from $-2$\textdegree C to
$-1$\textdegree C below normal raises growth rates in the cold scenarios
sufficiently that the same late-stocking plan remains executable, and
the structural infeasibility disappears. This is itself a useful confirmation that the 27/81 infeasibility ratio
in the main experiments is not an artefact of the mathematical
formulation but a genuine consequence of the cold tail of the scenario
tree breaking the EV-optimal stocking schedule.

EVPI drops substantially under reduced variance, from 128.1 to 73.8~MNOK
for the SP and from 97.7 to 56.3~MNOK for the DE, a 42\% reduction in
both cases. This is the expected behaviour: a narrower spread of futures
means perfect information has less to add, and the gap between WS and
the non-anticipative solutions shrinks accordingly. The VSS for the DE
follows the same pattern, falling from 57.0 to 50.0~MNOK as the cost of
ignoring uncertainty diminishes when the uncertainty itself is smaller.

The VSS for the SP appears to move in the opposite direction, rising from
26.6 to 32.5~MNOK, but this does not reflect a genuine increase in the
value of stochastic programming. Under the full-variance tree, EEV was
computed as a conditional expectation over only 54 feasible scenarios,
excluding the 27 where the EV plan failed entirely and EEV was therefore
artificially low. Under reduced variance all 81 scenarios are feasible
and EEV rises to 703.0~MNOK, which narrows the gap to SP from below.
The SP plan itself has not improved; the baseline it is measured against
has.

The objective gap between SP and DE remains tight at 2.3\% (735.5 versus
753.0 MNOK), and DE converges to a true 2\% optimum at this scale
(1.998\% reported gap) rather than terminating at the tolerance ceiling.
Wall-clock times are broadly comparable to the full-variance 10-cage run
(SP at 347 versus 386 seconds, DE at 1{,}946 versus 1{,}977 seconds),
indicating that the reduced spread does not materially change the
combinatorial structure of the model or the difficulty of the
binary-consensus problem the SP solves.


\subsection{Sensitivity to stages}
\label{sec:stages}

We next examine how VSS and EVPI respond to the granularity of the
scenario tree, holding fleet size at 60 cages and varying the number of
stages from two to four. The 4-stage configuration matches the main
experiments in §\ref{sec:sp_vs_ef}; the 2-stage and 3-stage variants
collapse the horizon into coarser temperature blocks while keeping the
3-branch (cold/normal/warm) structure at each stage, producing 9 and 27
scenarios respectively. Table~\ref{tab:model_stats_stages_60} reports
the result.

\begin{table}[htbp]
\centering
\caption{Model statistics and performance by stage count, 60 cages}
\label{tab:model_stats_stages_60}
\footnotesize
\begin{tabularx}{\linewidth}{lRRRRRRRRR}
\toprule
& \multicolumn{3}{c}{\textbf{2 Stages}} & \multicolumn{3}{c}{\textbf{3 Stages}} & \multicolumn{3}{c}{\textbf{4 Stages}} \\
\cmidrule(lr){2-4} \cmidrule(lr){5-7} \cmidrule(lr){8-10}
\textbf{Metric}
 & \textbf{EEV} & \textbf{SP} & \textbf{WS}
 & \textbf{EEV} & \textbf{SP} & \textbf{WS}
 & \textbf{EEV} & \textbf{SP} & \textbf{WS} \\
\midrule
Scenarios
 & 6/9\textsuperscript{*}   & 9  & 9
 & 18/27\textsuperscript{*} & 27 & 27
 & 54/81\textsuperscript{*} & 81 & 81 \\
E[obj] (MNOK)
 & 2{,}101 & 3{,}101 & 3{,}821
 & 2{,}854 & 3{,}306 & 4{,}041
 & 2{,}384 & 3{,}236 & 3{,}871 \\
MIP gap tolerance (\%)
 & 2 & 2 & 2
 & 2 & 2 & 2
 & 2 & 2 & 2 \\
Wall clock time (s)
 & 308 & 1{,}216 & 1{,}070
 & 1{,}076 & 2{,}110 & 2{,}429
 & 510 & 3{,}421 & 2{,}027 \\
\midrule
VSS$_{\text{SP}}$ (MNOK)
 & \multicolumn{3}{r}{$\text{SP} - \text{EEV} = 1{,}001$}
 & \multicolumn{3}{r}{$\text{SP} - \text{EEV} = 453$}
 & \multicolumn{3}{r}{$\text{SP} - \text{EEV} = 852$} \\
EVPI$_{\text{SP}}$ (MNOK)
 & \multicolumn{3}{r}{$\text{WS} - \text{SP} = 720$}
 & \multicolumn{3}{r}{$\text{WS} - \text{SP} = 734$}
 & \multicolumn{3}{r}{$\text{WS} - \text{SP} = 635$} \\
\bottomrule
\end{tabularx}
\par\vspace{2pt}
\footnotesize\textsuperscript{*}EEV expectation taken over the feasible scenarios only (probability mass 0.667 in all three columns).
\end{table}

A central finding is that the value of stochastic programming is
substantial at every stage count tested: VSS ranges from 453 to 1{,}001
MNOK and EVPI from 635 to 734 MNOK across the three configurations. The
SP captures a lot of the wait-and-see plan's advantage in every case. Ignoring uncertainty
consistently costs hundreds of MNOK regardless of how the horizon is
sliced, and the qualitative ordering EEV $<$ SP $<$ WS holds throughout.
SP wall-clock time also scales as expected: 1{,}216 s at 9 scenarios,
2{,}110 s at 27 scenarios, and 3{,}421 s at 81 scenarios, broadly
consistent with the linear-in-cages pattern documented in §\ref{sec:sp_vs_ef}. 

The non-monotonic VSS pattern (1{,}001 $\rightarrow$ 453 $\rightarrow$
852 MNOK) and the 14\% drop in EVPI between 3 and 4 stages is peculiar. Adding to the number of stages alters the
scenario tree in ways that go beyond simply improving the resolution. Even
though the expected temperature remains the same across all three
configurations, the underlying distribution of outcomes changes: the
number and length of stages directly determines how often two consecutive warm
realisations can occur and at which point in the horizon they fall.
Since elevated mortality in our model is triggered exclusively by two
consecutive warm stages, fewer stages means fewer opportunities for
high-mortality scenarios to occur, and the effective mortality risk
embedded in the scenario tree therefore differs across stage counts. 

The three stage-count columns
therefore describe genuinely different stochastic problems with
different effective uncertainty, not three resolutions of the same
problem, and comparing VSS and EVPI across them should be read with
that in mind. We also suspect that the 3-stage configuration may
interact favourably with the fallow structure of the instance: with 60
months divided into three stages of 20 months each, the non-anticipativity
constraints lock decisions over blocks that align closely with the
coordinated fallow periods of approximately 24 months, potentially
making it easier for the stochastic plan to coordinate stocking and
harvest decisions around those constraints than under the 15-month
blocks of the 4-stage tree. This also drives the unusual EEV pattern (2{,}101, 2{,}854, 2{,}384
MNOK), where the 3-stage EV plan happens to align more favourably with
its realised scenarios under that specific stage partitioning. A more
rigorous sensitivity study would fix the total horizon-level variance
and rescale the per-stage temperature deviations accordingly.

What survives the caveat is the qualitative finding the headline already
states: stochastic modelling earns its keep at every stage count tested,
and the stochastic solution captures a substantial fraction of the gap
to perfect information regardless of how the scenario tree is sliced.
The cross-stage comparison should nonetheless be read with the caveat
that the three trees represent different effective uncertainty, and a
controlled study rescaling the per-stage temperature deviations to hold
total horizon-level variance constant is left for future work.


\subsection{Sensitivity to Penalty Parameters}
\label{sec:sensitivity_params}

% """ BEHOLDE DETTE FØRSTE AVSNITTET??? """

% The PH algorithm is sensitive to the penalty weight $\rho$ and, more broadly, to
% the biomass-constraint penalty used in the objective.  Reducing the
% MAB slack penalty from 1{,}000~NOK/kg/month to 100~NOK/kg/month produced worse
% expected objectives, despite the looser penalty allowing more feasible solutions.
% This counter-intuitive result is consistent with the heuristic nature of PHA for
% MIPs: a weaker penalty changes the landscape of each iteration subproblem, leading the
% algorithm to different (and in this case inferior) consensus points.  It
% underscores that PHA does not guarantee global optimality for mixed-integer
% programs --- it is a high-quality heuristic whose solution quality depends on
% algorithmic hyperparameters.

\paragraph{Constraining regional MAB.} In practice, the regional Maximum Allowable Biomass (MAB) is often the most restrictive regulatory constraint for producers. However, in our initial dataset, this parameter was non-binding. To evaluate the model's sensitivity to this constraint, we conducted a sensitivity analysis across three levels of regional MAB using the 60-cage dataset.

The first instance is a \textbf{Relaxed} scenario with a limit of 35,000 tonnes. The second is a \textbf{Tangential} scenario set at 24,000 tonnes—just below the aggregate of all individual location MAB parameters: 
\begin{alignat*}{2}
    &= \text{loc}_1 + \text{loc}_2 + \text{loc}_3 + \text{loc}_4 + \text{loc}_5 + \text{loc}_6 \\
    &= 4,000\text{t} + 5,600\text{t} + 2,500\text{t} + 4,300\text{t} + 2,600\text{t} + 5,500\text{t} = 24,500\text{t}
\end{alignat*}
The third is a \textbf{Tight} scenario set at 18,000 tonnes. All models were run using the same experimental setup as in previous experiments, with performance metrics reported in Table \ref{tab:model_stats_MAB_diff}.

\begin{table}[htbp]
\centering
\caption{Model statistics and performance, 60 cages with different regional MAB limits}
\label{tab:model_stats_MAB_diff}
\footnotesize
\begin{tabularx}{\linewidth}{lRRRR}
\toprule
\textbf{Metric} & \textbf{Relaxed (35,000 tonnes)} & \textbf{Tangentially (24,000 tonnes)} & \textbf{Tight (18,000 tonnes)} \\
\midrule
Scenarios              & 81/81 feasible & 81/81 feasible & 81/81 feasible \\
E[obj] (MNOK)          & 3,235.5        & 3,233.6        & 3,089.9 \\
MIP gap tolerance (\%) & 2              & 2              & 2 \\
Wall clock time (s)    & 3,420.5        & 3,706.6        & 5,962.2 \\
\bottomrule
\end{tabularx}
\end{table}

As shown in Table \ref{tab:model_stats_MAB_diff}, tightening the regional MAB from 35,000 to 24,000 tonnes results in a negligible change to the expected objective value ($E[obj]$). This marginal difference is considerably insignificant given that our solution method is heuristic and utilizes a 2\% MIP gap, though it does demonstrate model robustness under moderate tightening. However, the computational effort already begins to scale, with wall-clock time increasing by approximately 8%

In the Tight scenario (18,000 tonnes), the expected profit decreases by roughly 145 million NOK (4.5\%). While this profit reduction is relatively modest compared to the 25\% reduction in allowable biomass, the computational cost is significant: wall-clock time increases by 60\% to 5,962 seconds. While this remains manageable for monthly tactical planning, it indicates a more complex solution space.

We hypothesize that this increased runtime is driven by the internal mechanics of the Gurobi solver’s \textbf{branch-and-cut} algorithm. When the regional MAB is set below the sum of location limits, it functions as a linking constraint. In combinatorial terms, while the regional MAB constraint technically strengthens the LP relaxation by tightening the dual bound, it significantly complicates the polyhedral structure of the problem as the constraint couples previously independent variables across different locations. This complicates the solver's ability to identify facet-defining inequalities that characterize the convex hull of the global feasible region. Instead of solving relatively independent sub-problems, the algorithm must navigate a more intricate search space where local cuts are less effective at pruning the search tree, forcing the solver to generate a significantly higher volume of cutting planes to reach the desired MIP gap.

\paragraph{Case study on unrealistic fallowing schedule}
While running the comparison between EEV and sp algorithm on 10, 30, and 60 cages we noticed the EEV was failing on a third of the scenarios, specifically, every scenario path starting with a cold stage. The reason was that location one, which was included in all the experiments, had its first scheduled fallowing month in month 12, $\mathcal{F}_1 = 12$. The EEV sees only a future in which growth occurs at normal temperatures. In this future, it manages to stock a cohort letting it grow past two kg and be eligible for harvest by the time it has its first scheduled location fallowing in month 12. Although our stochastic model foresees that under colder temperature scenarios this second cohort will not reach the two kg minimum the well-boat requires, the EEV does not, resulting in 27/81 scenarios being infeasible, as reported in the experiment section \ref{sec:sp_vs_ef}. The plan, under normal temperature realisation can be seen in figure \ref{fig:10_cage_plan}. 


This has mainly two interesting insights. Firstly, that the stochastic approach works as intended by planning more conservatively knowing that the future is not perfect. Secondly, that the parameters we feed our model and similarly the rules for production that is decided in the real world carries much importance. One might argue having $\mathcal{F}_1 = 12$ might be a unrealistic case and that we are in a sense setting the EEV up for failure under an unforeseen cold realisation with a tight fallowing schedule. If we set $\mathcal{F}_1 = 10$, then we end up with a more robust EEV feasible in 81/81 scenarios, as can be seen in figure \ref{fig:second}

% It is important to note that only showing model statistics where the model behaves in preferring manner is poor research which is why we decided to keep the original parameter setting for experiment sections \ref{sec:Determinisitic_compared_Stochastic} and \ref{sec:sp_vs_ef}. 

\begin{figure}[htbp] % [htbp] helps placement and prevents text "collision"
\centering
\begin{minipage}{0.48\textwidth}
    \centering
    \includegraphics[width=\linewidth]{bilder/sp_normal_path_timeline.png}
    \caption{This is plan the first 15 months with deterministic planning under normal temperatures with scheduled fallowing: $\mathcal{F}_1 = 12$. All cages are stocked at month zero and harvested at the latest by month 12.}
    \label{fig:10_cage_plan}
\end{minipage}
\hfill % Adds horizontal space between the two images
\begin{minipage}{0.48\textwidth}
    \centering
    \includegraphics[width=\linewidth]{bilder/10_cage_f10.png}
    \caption{First scheduled fallowing has $\mathcal{F}_1 = 10$. The model keeps away from stocking in month zero, and hence the bottom cages stays empty. The other lines are not stocked cohorts, but rather pre-existing cohorts that still reaches the two kg minimum requirement.}
    \label{fig:second}
\end{minipage}
\end{figure}

\newpage

% =============================================================
%  §5.9  Sensitivity to Horizon Length
% =============================================================
\subsection{Sensitivity to Horizon Length}
\label{horizon_lengths}
 
How far ahead the planner needs to look is a practical question with a compute cost: every additional month in the planning window widens the scenario tree and inflates solve time. We approach this in two steps. First, we isolate the effect of horizon length under a deterministic planner over a long horizon, where the comparison is clean and cheap. Second, we ask whether adding explicit scenario information changes the answer. To evaluate the performance of the proposed models in a dynamic environment, we employ a \textbf{Rolling Horizon Approach}. Rolling horizon is an iterative decision-making framework where the optimization model is solved over a multi-period planning horizon, but only the decisions for the immediate "roll" (or "commitment period") are executed. The horizon then shifts forward, the model state is updated with the realized data (in this case, realized temperature paths and salmon cohort status), and the process is repeated.
 
\paragraph{Deterministic horizon comparison.}
We roll the deterministic version of our model, solved only for normal/expected-values for temperature and mortality (EV-style MILP), forward across a 120-month horizon under two planning depths --- 30 and 60 months --- committing 12 months per roll for both models. Since the planner only assumes normal temperatures inside every roll, the plan structure is fixed and only the \emph{realised} temperatures vary. We evaluate performance across eight temperature paths, drawn uniformly for each stage from $\{\textsf{cold}, \textsf{normal}, \textsf{warm}\}$. Figure~\ref{fig:det_horizon_npv} overlays cumulative realised NPV per (variant, path). 

The 30-month variant averages 1.23 B NOK and outperforms the 60-month variant on colder and more variable paths (+16\%), while the 60-month variant wins on normal temperatures (+8\%), averaging 1.17 B NOK. Per-path differences range from $-0.10$~B~NOK to $+0.18$~B~NOK. The differences can stem from the 60-month plan making longer-reaching commitments under wrong assumption, so when temperature-anomalies hit, those commitments are harder to recover from in the next roll. The 30-month plan is wrong for a shorter period before it can replan. Difference may also be partly attributed to the solver's 3\% optimality gap, so we do not draw very strong conclusions from them. What the results does establish is: under full-information planning, there is no detectable mean-NPV advantage to either horizon over the 10-year real-time evaluation. The SP comparison below asks whether explicit uncertainty modelling changes that picture.

\begin{figure}[t]
  \centering
  \includegraphics[width=\linewidth]{bilder/deterministic_rolling_horizon_30_60.png}
  \caption{Cumulative realised NPV across a 120-month horizon under the deterministic plan-builder rolled forward under two planning depths. Faint lines are individual temperature paths; bold lines are per-variant means. Both planners assume all-normal weather inside every roll; the only plan-side difference between them is the per-roll terminal-value boundary at $t = 30$ versus $t = 60$ months.}
  \label{fig:det_horizon_npv}
\end{figure}

\paragraph{Stochastic horizon comparison on small dataset (30 cages).}
\label{par:horizon30}
Three rolling-horizon stochastic program variants were compared over a 120-month evaluation period: a 30-month horizon with 9 scenarios (\texttt{30M}), a 30-month horizon with 81 scenarios (\texttt{30M\_81}), and a 60-month horizon with 81 scenarios (\texttt{60M}). The (\texttt{30M}) and (\texttt{60M})variants still follow the 15 month stage lengths, while the (\texttt{30M\_81}) variant has these 15 months split into two stages of 8 and 7 months to create a scenario tree with 81 leaves, but still be comparable under the same realisations as the other two. The variants share the same realised temperature sequences, but differ in planning horizon and re-solve frequency: \texttt{30M} and \texttt{30M\_81} commit 15-month plans across eight re-solves, while \texttt{60M} commits 45-month plans across three.

The \texttt{60M} variant delivers substantially higher NPV across every temperature trajectory (Figures~\ref{fig:npv} and~\ref{fig:npvbars}), with a feasibility-weighted average of 1\,486 MNOK, against 726 MNOK for \texttt{30M}. Importantly, this advantage emerges despite a structural disadvantage: by committing to 45-month plans, \texttt{60M} updates its plan only three times over the evaluation horizon, less than half as frequently as the 30-month variants. That the longer horizon nonetheless dominates by such a margin points to the value of foresight in this setting. We argue that for this setup, the ability to account for future biomass accumulation and weight-class progression across several production cycles appears to outweigh the cost of reduced responsiveness to new temperature information.


Introducing richer scenario structure within the same 30-month horizon \texttt{30M\_81} yields mixed results. On normal or consistently cold trajectories, \texttt{30M\_81} modestly outperforms \texttt{30M\_81}, suggesting that finer uncertainty branching can improve decision making for some outcomes. On warm or stress-then-recover trajectories however, it performs worse. A likely cause one can imagine is the hedging cost: stage decisions covering only seven months must remain feasible across a much more uncertain sample space, and the resulting conservatism is costly when the actual trajectories have much less variation. For these temperature variants, expanding the scenario tree is (obviously enough) not beneficial unless the planning horizon is extended alongside it. If we were to compare them under finer-grained temperature realisations, varying across 7-8-month stages rather than across 15-month ones, then it might better play to the \texttt{30M\_81} plan's advantage, and it could alter the relative performance seen here. This touches on the same type of comparison issues raised in section \ref{sec:stages}. 

The computational cost of rolls per trajectory for the three variants differs by an order of magnitude (Figure~\ref{fig:clock}). The lower re-solve frequency of \texttt{60M} and the shorter horizon of the \texttt{30M\_81} partially offsets this, but overall wall time is substantially longer than that of \texttt{30M}. In an operational context of a larger fleet size or a more detailed uncertainty representation, this gap would be non-trivial and may constrain the practical applicability of longer-horizon models.

For producers operating a setup similar to the one modelled here, longer planning horizons appear to deliver meaningful economic value, possibly enough to justify the added computational cost. An open question is whether a \texttt{60M} horizon with a shorter commit window --- re-solving every 15 months rather than every 45 --- would yield further gains, at greater computational cost, or whether the current results already reflect a reasonable balance between foresight and adaptability.

\begin{figure}
    \centering
    \includegraphics[width=1\linewidth]{bilder/trajectories.png}
    \caption{The five different temperature realisations we asses the different model variants on.}
    \label{fig:temp_traj}
\end{figure}

\begin{figure}
    \centering
    \includegraphics[width=1\linewidth]{bilder/30v60trajprofit.png}
    \caption{The different net present valued profits of the 120 month realisations for the different variants' plan. The \texttt{30M\_81} variant got stuck solving the initial per scenario problems on the fifth roll under oscillating temperatures and is therefore omitted.}
    \label{fig:npv}
\end{figure}

\begin{figure}
    \centering
    \includegraphics[width=1\linewidth]{bilder/profit_npv.png}
    \caption{The accumulated profit over the 120 month long horizon discounted back to month zero for the different variants and temperature trajectories/realisations.}
    \label{fig:npvbars}
\end{figure}

\begin{figure}
    \centering
    \includegraphics[width=1\linewidth]{bilder/wall clock horizon exp.png}
    \caption{The time (in hours) it took to run each full trajectory including all horizon rolls.The oscillating run for \texttt{30M\_81} is omitted her as the solver got stuck on roll 5.}
    \label{fig:clock}
\end{figure}
 
\paragraph{Stochastic horizon comparison on large dataset (60 cages) with deterministic tale.} To fully analyse the horizon effects we also take a look at the behaviour on a more realistic dataset size of a 60 cage regional fleet and run two SP horizon variants --- a 30-month branching window and an extended 60-month plan. Unlike the deterministic comparison above, which commits either 15 or 45 months per roll, the stochastic variants here commit half their planning horizon per roll. A structural note on the 60-month variant is necessary before reading the results. Its scenario tree is \emph{identical} to the 30-month SP: a 3-month deterministic prefix followed by three 9-month stages branching on $\{\textsf{cold}, \textsf{normal}, \textsf{warm}\}$, giving $|\mathcal{S}| = 27$ scenarios. The additional 30 months are a \emph{deterministic EV continuation}: a single all-normal temperatures tail appended to all 27 scenarios. The 60M variant therefore does not represent 60 months of stochastic information. The comparison is more precisely: does extending the planning horizon beyond the branching window, via an EV continuation tail, affect the committed first-stage decisions or realised profit to a large degree? This experiment setup is a tractability move to deal with repeated long-horizon problems, not to be confused with richer scenario information.
 
Both runs use the MIP gap tolerance of 3\%. We evaluate them under three in-sample temperature paths, each composed of two 30-month regimes at the (prefix $= 3$, $s_1 = 9$, $s_2 = 9$, $s_3 = 9$) block granularity:
 
\begin{figure}
    \centering
    \includegraphics[width=1\linewidth]{bilder/30_60_profit_longterm.png}
    \caption{Cumulative realised NPV over the 60-month real horizon under different temperature paths. The left plot shows run under normal temperatures, the middle under mixed temperatures (first half normal, second half oscillating), and the right under oscillating temperatures. Blue solid: 30-month SP. Red dashed: 60-month SP with deterministic EV continuation tail.}
    \label{fig:sp_horizon_npv}
\end{figure}
 
% \begin{figure}
%     \centering
%     \includegraphics[width=1\linewidth]{bilder/03_per_solve_wallclock.png}
%     \caption{Enter Caption}
%     \label{fig:wallclock}
% \end{figure}

Three paths is too small a sample for inference, and the comparison carries structural asymmetries that limit interpretation: the two variants commit different absolute lengths of real time per roll, and the 30M variant may have a relative advantage because month 60 is the end of its planning horizon, while the 60M variant is only measured halfway through its horizon on the second roll. We therefore limit ourselves to a single observation the data does support: across all three paths, the 60M variant, which only adds a deterministic extended horizon, has a substantially higher per-solve wallclock time without producing a consistent profit improvement within the evaluation window. A longer \textbf{stochastic} horizon appears to deliver more meaningful value, as the preceding experiment demonstrates.

\newpage

% =============================================================
%  §5.9  An Illustrative Implementation of the Rolling SP
% =============================================================
\subsection{An Illustrative Implementation of the Rolling SP}
\label{ssub:rh_experiment}

 
The following simulation experiment runs the SP-model framework iteratively through time. The purpose is to demonstrate plausible operational behaviour. It is illustrative rather than benchmark-grade.
 
\paragraph{Setup.}
Each roll plans for $T = 30$ months and commits the first three. The planning window consists of a 3-month deterministic prefix (the temperatures already realised) followed by three nine-month stages, each branching on $\{\textsf{cold}, \textsf{normal}, \textsf{warm}\}$, giving $|\mathcal{S}| = 27$ scenarios per roll. A shorter horizon and fewer scenarios is chosen to shorten the simulation time duration. After PH terminates, the consensus first-stage decisions are committed and the model state --- biomass per cohort, location occupancy, regional MAB usage --- advances under the realised temperatures before the next roll re-solves from the updated state. Twelve rolls cover a 36-month real horizon against the five temperature realisations shown in figure \ref{fig:temp_traj}, just adapted to serve nine-month long stages instead of 15, and the \textit{cooling} and \textit{warming} have periods of normal temperatures, as can be seen in figure \ref{fig:warming}'s reality path. All 60 PH solves complete with sub-MILP gap at 3\% in approximately 5.3~hours.

\begin{figure}
    \centering
    \includegraphics[width=1\linewidth]{bilder/warming_sim.png}
    \caption{An example realisation of production under the \textit{warming} trajectory over this three year simulation rolling the horizon and re-planning every three months. The top graph shows which temperatures were realised during which months (both normal and warm temperatures). The bottom plot show how the standing biomass (blue line) and harvested biomass (yellow bars) progresses over the timeline realisation.}
    \label{fig:warming}
\end{figure}

\begin{figure}[t]
  \centering
  \includegraphics[width=\linewidth]{bilder/biomass_overlay.png}
  \caption{Realised standing biomass across the five temperature paths, with the regional MAB cap overlaid.}
  \label{fig:rh_overlay}
\end{figure}
 
\paragraph{Operational behaviour.}
Figure~\ref{fig:rh_overlay} overlays realised standing biomass under the five paths against the regional MAB cap. Under \textit{warming} scenarios, harvest timing advances by a couple of months relative to the \textsf{all\_normal} baseline, possibly clearing biomass ahead of predicted heat episodes and elevated mortality rates. The paths largely track each other through the first year ---while cohort sizes differ across temperature realisations, the initial stocking decisions were made under the same information --- and begin to diverge from around month~23 onward as the accumulated temperature differences start to affect growth rates and harvest windows. 

\begin{figure}
    \centering
    \includegraphics[width=1\linewidth]{bilder/fanchart_cold.png}
    \caption{Per-roll biomass fan charts for the \textsf{cooling} temperature path (12 panels, one per roll). Each panel shows all 27 scenario trajectories from the final PH iteration of that roll, together with the probability-weighted mean (black) and the the commit point, three months further into each new horizon (red vertical line).}
    \label{fig:fanchart}
\end{figure}

Figure~\ref{fig:fanchart} shows the per-roll fan charts for the \textsf{cooling} path as a representative for the five trajectories. The most notable pattern is how the plan revises itself as the horizon advances. The model starts of by planning a substantial mid-horizon biomass peak (around month 12). Interestingly, by the time those months are actually reached, the model has prolonged the biomass increase in the short term and the dip that appeared in the early plans does not materialise in the committed decisions until about a year later. This is perhaps an illustration of the rolling framework: decisions that look overly conservative when taken early are revised continuously as uncertainty resolves, and the plan at any point in time reflects only the uncertainty that genuinely remains.

\newpage

\section{Conclusion}
\label{sec:Aftermath}

\subsection{Modelling the Risk}
Through the thesis we have only been looking at one metric: maximizing expected profit. This is a typical way of looking at a production problem as most producers actually do want as much profit as possible. It is also the challenge our collaborators at Scale AQ initially gave us, namely maximize the total biomass. However, äll this treats the optimizer as risk neutral which is not always the case. Some producers might be willing lower their expected profit for a higher minimum return or less variance in the profit outcomes. Without going into much detail, we can disclaim that there are more options than what we have used. One example, is the Conditional Value at Risk (CVaR) method, maximizing the average profit across the x\% worst scenarios instead. The disadvantage with this is the model's potential to "sacrifice" a high expected value plan for a plan with a fractionally better downside. The Mean-CVaR method, which optimises over a weighted sum of expected profit and the x\% worst cases, might also be a good alternative. There are a lot of ways one could choose to model the risk-factor and the point is that producers have different priorities when it comes to capturing the right amount and type of risk. We believe this model-algorithm framework has the potential to be adapted to anyone's risk-preferences by altering the objective function of the model. 

% =============================================================
%  §5.11  Suggestive Implementation
% =============================================================
\subsection{Suggestive Implementation}
\label{sec:impl-recommendation}
 
% NOTE: This section is a work in progress. Specific recommendations
% on horizon length and MIP gap depend on the stage-sensitivity and
% scalability experiments; update once those results are final.
 
A operational adaptation of the SP algorithm for salmon producers is a frequent rolling horizon re-solve: new temperature observations can be adjusted monthly, so committing one month at a time and re-solving from the updated state keeps the plan maximally responsive to new information. Our experiments show that a longer planning horizon and a deeper scenario tree are beneficial. However, they are not independent choices, but rather jointly constrained. For a small number of cages one may afford a longer planning horizon and a richer scenario tree. For a non-binding regional MAB, the production fleet can be divided into much simpler sub-problems. For larger fleets on the  other hand, the computational time dominates and a sparser/shallower tree or a shorter timeline may be necessary to stay within budget. The same applies to the MIP gap tolerance: a tighter gap increases the quality of each scenario's solution but raises per-iteration cost. The right balance for all these modelling choices depends on where the model convergence bottleneck lies and where the most gain can be made for a given instance. 

 
% \subsection{Discussion - What could have been done better?}
% Modeling the true temperature distribution, and subsequenbtly the true growth and mortality. The arguably too large temperature difference between the cold, normal and warm realisations can also be seen as a way to envelop some of the stochasticity in salmon farming we could not capture like fish escapes, production hiccups and other uncertainties. If this approach is positive or negative for production planning is something we would gladly be proven at.

% Our model does not offer a full recourse if the plan ends up infeasible for certain scenarios, and the EEV is not measurable under certain unforeseen scenarios even if they are not that harsh.
% % """ Diskutere rundt hvor viktig det er å modellere usikkerheten riktig vs. det å bare modellere den i det hele tatt. Spør Stein Ove om artikkelen og forskeren som så på nettop dette på NTNU..."""

% \subsection{Future Research}
% More Realism:
% - Incorporating price uncertainty and seasonal supply smoothing incentive.
% - More harvest opportunities per cohort (harvest half now, half later)
% - More smolt sizes

% VSS and out of sample testing:
% - See how well our stochastic model vs deterministic model responds to historical temperature data instead of the "reality" we built. This is hard to do as we have not modelled this for a plan that responds to intervals of temperature realizations. 

\subsection{Summary}
\label{sec:conclusion}

This thesis has developed and evaluated a stochastic mixed-integer linear programming approach to long-term production planning in Atlantic salmon farming. Starting from the deterministic MILP of \citet{ErikstadHaugland2025} and the operational framing supplied by Scale Aquaculture, we formulated a four-stage SMILP with 81 scenarios capturing temperature-driven uncertainty in growth and mortality, and proposed a tailored Progressive Hedging matheuristic to solve it at industrial scale.

The central modelling idea is to separate the combinatorial axis of the problem from the recourse axis. Stocking and harvest timing are binary, and their non-anticipativity is what couples scenarios into an intractable optimisation. Stocking quantities, by contrast, are continuous in practice and well behaved once the timing skeleton is fixed. Rather than enforcing consensus on all non-anticipative variables simultaneously, the algorithm drives only the binary decisions to agreement through Lagrangian penalties and progressive variable fixing, then solves a single extensive-form stochastic LP in the stocking quantities to find the optimal stocking decisions given the committed binary plan. 
%This design reflects a structural feature of the salmon production planning problem rather than a generic decomposition trick. 
The consensus space for binaries is finite and combinatorially small, while consensus on continuous quantities like $q$ would require agreement on real-valued numbers in the tens of thousands and is what makes Progressive Hedging notoriously slow to settle on mixed-integer problems. The two simplifications compound: the PH phase converges on a much smaller search space than a full-variable PH would face, and the recourse phase reduces to an extensive-form stochastic LP in which non-anticipativity on $q$ can be enforced directly, since the problem is now convex and continuous.

The empirical results support three conclusions. First, on the instances where the deterministic equivalent is still tractable, the B-PHA heuristic creates solutions within roughly five percent of the DE objective, and the gap is consistent with the combined MIP gap budget of the two solvers. The heuristic is therefore sacrificing some solution quality, although not too much so, at scales where the comparison can be made directly. Second, the heuristic scales approximately linearly in the number of cages while DE wall-clock time grows by more than an order of magnitude per doubling of fleet size and exhausts memory entirely at 40 cages on the test hardware. The 17-fold speedup at 30 cages and the qualitative tractability at 60 and 70 cages establish B-PHA as the only approach in the comparison that is both reliable and applicable at industrial fleet sizes on commodity hardware. Third, biological uncertainty has substantial economic value in this setting. The deterministic expected-value plan is infeasible in 27 of 81 scenarios at every fleet size, a robustness failure that the conditional VSS metric understates by construction. Both VSS and EVPI grow with fleet size, reflecting the accumulation of independent committed errors that a deterministic plan cannot hedge, and both shrink substantially under a narrower temperature spread, consistent with the principle that stochastic planning earns its keep in regimes where realisations differ substantially from each other.

\subsection{Limitations}

Several limitations bound these conclusions. The scenario tree discretises continuous temperature variables into three branches per stage, and consequently, the reported VSS and EVPI are conditional on this construction rather than on the empirical distribution of Norwegian sea temperature or an out-of-sample benchmark. We model temperature as the only stochastic source and treat price as a deterministic risk-adjusted parameter; in practice salmon spot prices, institutional changes and fish diseases are all sources of uncertainties, and \citet{bergfjord2009} reports that producers themselves rank these as their dominant source of risk. The independence assumption across stages is also a simplification, and explicitly correlated temperature realisations would make the scenario tree more realistic, although at the cost of additional modelling complexity. On the algorithmic side, B-PHA is a heuristic without convergence guarantees, and its reliability rests on empirical observation that the binary-only consensus problem is well behaved for this problem
instances rather than on a formal optimality proof.


\subsection{Further Work}
Numerous directions follow from this work. The most direct extension is to incorporate price uncertainty alongside biological uncertainty. Another extension would be to make prices depend on harvesting month, addressing seasonal fluctuations and the potential benefits of smoothing harvest throughout the year. Furthermore, a more refined biological recourse, such as partial harvest of a cohort or an emergency harvest option, would expand the recourse structure and likely increase the value of stochastic planning. Integrating multiple smolt sizes is also a logical next step. Although we initially explored this, it was deferred due to the combinatorially explosive impact on the size of the problem. On the algorithmic side, a decrease in memory footprint or increase in the memory budget could push the empirical ceiling beyond the 80-cage limit observed here, and out-of-sample stability testing on historical temperature data would speak more directly to the realised value of stochastic planning than the in-sample VSS reported in this thesis.

\subsection{Concluding Remarks}
The broader takeaway is that long-term salmon production planning under biological uncertainty is computationally tractable at scales relevant for industrial production levels, and that the gain from doing so is large enough to be worth the modelling effort. The true value of the stochastic approach lies in its ability to reconcile economic optimization with operational reality. While deterministic models frequently produce plans that become unexecutable when faced with biological volatility, the B-PHA heuristic ensures feasibility across a wide range of scenarios without sacrificing competitive expected NPV. By maintaining this reliability at a computational cost that scales modestly, the model demonstrates that sophisticated stochastic planning is a viable alternative to traditional deterministic frameworks in the face of an increasingly uncertain biological environment. 
%The value of stochastic planning is not primarily about improved expected profit on plans that would be feasible anyway; it is about producing plans that remain executable across the range of futures the producer actually faces. On that criterion, the deterministic baseline fails in roughly a third of scenarios at every fleet size considered, while the B-PHA heuristic remains executable in all of them at a wall-clock cost that grows far more slowly than the extensive form and only modestly above the deterministic plan itself. Under low-variance uncertainty it still delivers measurable expected-NPV gains; under high-variance uncertainty it secures feasibility and competitive expected NPV where the deterministic alternative delivers neither.
\newpage

\bibliographystyle{plainnat}
\bibliography{references}

% \section{Appendix}


\end{document}